% concatenated from paper.tex, header.tex, appendix.tex
%%%%%%%% ICML 2026 EXAMPLE LATEX SUBMISSION FILE %%%%%%%%%%%%%%%%%
\documentclass{article}

% Use the following line for the initial blind version submitted for review:
% \usepackage{icml2026}

% For preprint, use
% \usepackage[preprint]{icml2026}

% If accepted, instead use the following line for the camera-ready submission:
\usepackage[accepted]{icml2026}

% Recommended, but optional, packages for figures and better typesetting:
\usepackage{microtype}
\usepackage{graphicx}
% \usepackage{subfigure}
\usepackage{booktabs} % for professional tables

% hyperref makes hyperlinks in the resulting PDF.
% If your build breaks (sometimes temporarily if a hyperlink spans a page)
% please comment out the following usepackage line and replace
% \usepackage{icml2025} with \usepackage[nohyperref]{icml2025} above.
% \usepackage{hyperref}

% Attempt to make hyperref and algorithmic work together better:
% \newcommand{\theHalgorithm}{\arabic{algorithm}}

%%%%%%%%%%%%%%%%%%%%%%%%%%%%%%%%%%%%%%%%%%%%%%%
%%%% OWN COMMANDS
%%%%%%%%%%%%%%%%%%%%%%%%%%%%%%%%%%%%%%%%%%%%%%%

%%%%%%%%%%%%%%%%%%%%%%%%
% IMPORTS
%%%%%%%%%%%%%%%%%%%%%%%%

\usepackage{xcolor}
\usepackage{amsmath,amssymb,amsfonts}
\usepackage{mathtools}
\usepackage{amsthm}
\usepackage{bm}
\usepackage{dsfont}
\usepackage[]{algorithmic, algorithm}
\usepackage{enumitem}
\usepackage{natbib}
\usepackage{verbatim}

%%%%%%%%%%%%%%%%%%%%%%%%%%%%%%%%%%%%%%%%%%%%%%%%%%%%%%%%%%%
%%%%%%%%%%%%%%%%%%%%%%%%%%%%%%%%%%%%%%%%%%%%%%%%%%%%%%%%%%%

% create tables with multi-rows
\usepackage{multirow}
\usepackage{longtable}

% for links in pdf, citing urls
\usepackage[hidelinks]{hyperref}
\usepackage[capitalize,nameinlink]{cleveref}
% NOTE: thmtools makes link go to the wrong location
% \usepackage{thmtools}
% \usepackage{thm-restate}

\usepackage{subcaption} % for subfigures
\usepackage{caption}

\usepackage{url}
\usepackage{cite}

%%%%%%%%%%%%%%%%%%%%%%%%
% COMMANDS
%%%%%%%%%%%%%%%%%%%%%%%%

%switch vareps and varphi
\let\temp\phi
\let\phi\varphi
\let\varphi\temp

\let\temp\varepsilon
\let\epsilon\varepsilon
\let\varepsilon\temp

% optimization and convex analysis
\DeclareMathOperator*{\argmin}{arg\,min}
\DeclareMathOperator*{\argmax}{arg\,max}
\newcommand{\prox}[1]{\mathrm{prox}_{#1}}
\newcommand{\env}[1]{\mathrm{env}_{#1}}
\DeclareMathOperator{\conv}{conv}

% standard maths stuff
\newcommand{\iprod}[2]{\langle #1, #2 \rangle}

\newcommand{\then}{\Longrightarrow}
\newcommand{\onehalf}{\frac{1}{2}}
\newcommand{\oneover}[1]{\frac{1}{#1}}
\newcommand{\toneover}[1]{\tfrac{1}{#1}}
\newcommand{\T}{\top}

\newcommand{\R}{\mathbb{R}}
\newcommand{\Rext}{\overline{\mathbb{R}}}
\newcommand{\Rn}{\mathbb{R}^n}
\newcommand{\N}{\mathbb{N}}

\newcommand{\E}{\mathbb{E}}
\newcommand{\Prob}{\mathbb{P}}

\newcommand{\evmin}[1]{\lambda_{\min}(#1)} % min eigenvalue
\newcommand{\evmax}[1]{\lambda_{\max}(#1)} % max eigenvalue
\newcommand{\symm}{\mathbb{S}} % symm matrices
\newcommand{\SPSD}{\mathbb{S}_{+}} % symm pos semidefdef matrices
\newcommand{\SPD}{\mathbb{S}_{++}} % symm pos def matrices


\DeclareMathOperator{\Tr}{tr} % trace of matrix
\DeclareMathOperator{\dist}{dist}
\DeclareMathOperator{\dom}{dom}
\DeclareMathOperator{\diag}{diag} % extracts diagonal from matrix
\DeclareMathOperator{\Diag}{Diag} %creates matrix from vector
\DeclareMathOperator{\ri}{ri} % relative interior
\DeclareMathOperator{\intr}{int} % relative interior

\newcommand{\unitvec}[1]{\bm{\mathrm{e}}_{#1}} % unit vectors
\newcommand{\ones}[1]{\bm{1}_{#1}} % vector of ones
\newcommand{\indfun}[1]{\mathds{1}_{#1}} % indicator function 0/1
\newcommand{\elemfun}[1]{\bm{\delta}_{#1}} % indicator function 0/1
\newcommand{\Id}{\bm{\mathrm{Id}}} % identity
\newcommand{\Proj}[1]{\mathrm{Proj}_{#1}} % projection
\DeclareMathOperator{\sign}{sign} % sign

% bold math letters
\newcommand{\bw}{\bm{w}}
\newcommand{\be}{\bm{e}}
\newcommand{\bd}{\bm{d}}
\newcommand{\bg}{\bm{g}}
\newcommand{\bx}{\bm{x}}
\newcommand{\by}{\bm{y}}

% algorithms
\newcommand{\Adam}{{\texttt{Adam}}}
\newcommand{\SGD}{{\texttt{SGD}}}
\newcommand{\SGDM}{{\texttt{SGD-M}}}
\newcommand{\SPP}{{\texttt{SPP}}}
\newcommand{\SPS}{{\texttt{SPS}}}
\newcommand{\NGN}{{\texttt{NGN}}}

% plain text 
\newcommand{\ie}{i.e.\ }
%\newcommand{\skipline}{\vspace{\baselineskip}\\}


%%%%%%%%%%%%%%%%%%%%%%%%
%%% STYLING
%%%%%%%%%%%%%%%%%%%%%%%%
\definecolor{kleinblue}{RGB}{0, 47, 167}
\definecolor{royalblue}{RGB}{0, 33, 115}
\hypersetup{
	colorlinks = True,
	linkcolor = royalblue,
	%linkcolor = kleinblue,
	citecolor=royalblue,
	urlcolor=royalblue
} 

%%%%%%%%%%%%%%%%%%%%%%%%
%%% USEFUL
%%%%%%%%%%%%%%%%%%%%%%%%
\definecolor{todored}{RGB}{189, 30, 30}
\newcommand\revise[1]{{\color{todored}{#1}}}
\definecolor{nypink}{RGB}{216, 131, 115}
%\newcommand\todo[1]{{\color{nypink}{#1}}}
\definecolor{commblue}{RGB}{0, 66, 102}
\newcommand\comm[1]{{\color{commblue}{#1}}}

\definecolor{takeawaycolor}{RGB}{155, 157, 137}
\definecolor{StrongRed}{RGB}{230, 57, 70}
\definecolor{OtherBlue}{RGB}{26, 101, 158}
\definecolor{DarkGreen}{RGB}{39,130,67}
\definecolor{DarkRed}{RGB}{192,47,25}


% Orange boxes for comments
\usepackage[colorinlistoftodos]{todonotes}

\newcommand{\fabian}[1]{\todo[inline,bordercolor=orange,backgroundcolor=orange!20,linecolor=orange]{\textbf{Fabian: }#1}}
\newcommand{\takeaway}[1]{\todo[inline,bordercolor=takeawaycolor,backgroundcolor=takeawaycolor!20,linecolor=takeawaycolor]{\textbf{Takeaway: }#1}}

%%%%%%%%%%%%%%%%%%%%%%%%


%% !!!!!!!!!!!!!! NOTE !!!!!!!!!!!!!!!!!!!!!
% in the 2025 compiler of Latex, this setup always print crefs to Lemmas as Theorem XX.X
% we fix this here by hardcoding \label[lemma]{} everytime
% not sure how to fix this elegantly

%%%%%%%%%%%%%%%%%%%%%%%%%%%%%%%%
% THEOREMS
%%%%%%%%%%%%%%%%%%%%%%%%%%%%%%%%
\theoremstyle{plain}
\newtheorem{theorem}{Theorem}[section]
\newtheorem{proposition}[theorem]{Proposition}
\newtheorem{lemma}[theorem]{Lemma}
\newtheorem{corollary}[theorem]{Corollary}
\theoremstyle{definition}
\newtheorem{definition}[theorem]{Definition}
\newtheorem{assumption}[theorem]{Assumption}
\theoremstyle{remark}
\newtheorem{remark}[theorem]{Remark}


% change texttt font
\renewcommand{\ttdefault}{cmtt}
\usepackage[T1]{fontenc}


% The \icmltitle you define below is probably too long as a header.
% Therefore, a short form for the running title is supplied here:
\icmltitlerunning{Step-Size Stability in Stochastic Optimization}

\begin{document}

\twocolumn[
  \icmltitle{Step-Size Stability in Stochastic Optimization: A Theoretical Perspective}

  % It is OKAY to include author information, even for blind submissions: the
  % style file will automatically remove it for you unless you've provided
  % the [accepted] option to the icml2026 package.

  % List of affiliations: The first argument should be a (short) identifier you
  % will use later to specify author affiliations Academic affiliations
  % should list Department, University, City, Region, Country Industry
  % affiliations should list Company, City, Region, Country

  % You can specify symbols, otherwise they are numbered in order. Ideally, you
  % should not use this facility. Affiliations will be numbered in order of
  % appearance and this is the preferred way.
  \icmlsetsymbol{equal}{*}

  \begin{icmlauthorlist}
    \icmlauthor{Fabian Schaipp}{inria}
    \icmlauthor{Robert M.\ Gower}{fi}
    \icmlauthor{Adrien Taylor}{inria}
  \end{icmlauthorlist}

  \icmlaffiliation{inria}{Inria, Departement d'Informatique de l'Ecole Normale Superieure, PSL Research University, Paris, France}
  \icmlaffiliation{fi}{CCM, Flatiron Institute, New York City}
  %\icmlaffiliation{sch}{School of ZZZ, Institute of WWW, Location, Country}

  \icmlcorrespondingauthor{Fabian Schaipp}{fabian.schaipp@inria.fr}

  % You may provide any keywords that you find helpful for describing your
  % paper; these are used to populate the "keywords" metadata in the PDF but
  % will not be shown in the document
  \icmlkeywords{Optimization, Polyak step size}

  \vskip 0.3in
]

% this must go after the closing bracket ] following \twocolumn[ ...

% This command actually creates the footnote in the first column listing the
% affiliations and the copyright notice. The command takes one argument, which
% is text to display at the start of the footnote. The \icmlEqualContribution
% command is standard text for equal contribution. Remove it (just {}) if you
% do not need this facility.

% Use ONE of the following lines. DO NOT remove the command.
% If you have no special notice, KEEP empty braces:
\printAffiliationsAndNotice{}  % no special notice (required even if empty)
% Or, if applicable, use the standard equal contribution text:
% \printAffiliationsAndNotice{\icmlEqualContribution}

\begin{abstract}
  We present a theoretical analysis of stochastic optimization methods in terms of their sensitivity with respect to the step size.
    We identify a key quantity that, for each method, describes how the performance degrades as the step size becomes too large.
    For convex problems, we show that this quantity directly impacts the suboptimality bound of the method.
    Most importantly, our analysis provides direct \emph{theoretical} evidence that adaptive step-size methods, such as \SPS{} or \NGN{}, are more robust than \SGD{}.
    This allows us to quantify the advantage of these adaptive methods beyond empirical evaluation.
    Finally, we show through experiments that our theoretical bound qualitatively mirrors the actual performance as a function of the step size, even for non-convex problems.
\end{abstract}

%%
\section{Introduction}
%%
%
Training machine learning models requires algorithms that are at the same time efficient and robust. Several studies have shed light on the effect of subtle choices in architecture, loss function or optimizer on training efficiency and stability \citep{Wortsman2024,Everett2024,Olmo2025}. Specifically \citet{Wortsman2024} focus on the effect on the range of learning rates, for which training proceeds smoothly. This is not surprising, as it is widely accepted in the optimization and machine learning community that the learning rate is arguably the most important hyperparameter to tune \citep{Bengio2012}.

For this paper, we will focus on the effect of the \emph{optimization algorithm} on training stability. In the past, several methods have been designed to facilitate or even fully avoid the issue of tuning the learning rate (more in \cref{sec:related-work} below). While empirical evaluation in machine learning is closely paying attention to the issue of learning-rate sensitivity, optimization theory on the other hand puts more focus on convergence rates (where it is sufficient that the optimal rate is attained for \emph{some} step size).

To fill in this gap, here we compare -- with a theoretical framework -- how robust several \emph{stochastic} optimization methods are with respect to their step size. In particular, we focus on stochastic gradient descent (\SGD{}), the stochastic Polyak step size (\SPS{}), non-negative Gauss-Newton steps~(\NGN{}), and stochastic proximal point (\SPP{}). Throughout the paper we will use the term \emph{sensitivity} to refer to the final loss/suboptimality as a function of (the scale of) the step size: for small step sizes, all methods will only progress slowly; for large step sizes however, some methods behave more stably than others.

We summarize our contributions below.

\paragraph{Summary and contributions.} 
\begin{itemize}
    \item We identify a key quantity within our theoretical framework, which we call the \emph{stability index}, for determining how the method's suboptimality behaves as a function of the step size. Our theory covers a general family of model-based methods and applies to \SGD{}, \SPS{} \SPP{}, and partially \NGN{} in the convex and weakly convex setting. 
    \item Previous \emph{empirical} results suggest that \SPS{}, \NGN{}, and \SPP{} are always less sensitive to large step sizes compared to \SGD{}.  We will show that this is provably true.
    In particular, our results show that \textbf{\NGN{} and \SPS{} are always at least as stable as \SGD{}}. For \SPS{}, this holds for any guess of the optimal value/lower bound; interpolation is not necessary to obtain an advantage, but it helps.
    Our analysis also reveals the key quantities to determine the gain in stability, which are related to interpolation and the under-estimation of the optimal (batch) loss. 
    \item In several convex and non-convex toy experiments (linear/logistic regression, training a \texttt{ResNet} on \texttt{CIFAR10}), we track the theoretical quantities previously obtained to explain the improved stability of \SPS{}, \NGN{}, and \SPP{}. We show that the actual performance and the theoretical suboptimality bound match closely in predicting the range of learning rates that give good performance. Surprisingly, this is also true for non-convex problems, where our theory is not applicable.
\end{itemize}

\paragraph{Limitations.} The goal of this paper is predominantly to explain empirically observed phenomena with a general theoretical framework. Even though our experiments suggest a close match between theory and practice, we leave it open how to exploit this for practical applications. Further, some parts of our theory are restricted to convex problems, and only to a specific family of methods; in particular, it does not cover momentum, or preconditioning methods such as \Adam{}, which are prevalent in practice.


%
\paragraph{Setup and background.}
%
We consider the training problem
\begin{align}\label{eqn:problem}
    \min_{x\in \R^d} f(x), \quad f(x) :=  \E_{s}[f(x,s)].
\end{align}
Here, $x \in \R^d$ are the learnable parameters of a machine learning model, and $f$ is the loss function. In $f$, the expectation is taken over the distribution of a random variable $s$ that maps to the space or set $\mathcal{S}$ (typically the training set).
We assume that $f(\cdot,s)$ has a suitable subdifferential for every $s\in \mathcal{S}$ (for example, see \citet{Rockafellar1970,Clarke1983,Bolte2021}). We denote elements of the subdifferential by $g\in \partial f(x,s)$. We assume that problem \eqref{eqn:problem} admits a solution $x_\star$ and denote $f_\star := f(x_\star)$.

We briefly introduce the four methods we will mainly focus on throughout: for a step size $\alpha>0$\footnote{For now we assume a constant step size purely to keep the presentation simple; all results later work for time-dependent step sizes $\alpha_t$.}, \emph{stochastic gradient descent} is given by
\begin{align}\label{eqn:sgd}\tag{\SGD{}}
    x_{t+1} = x_t - \alpha g_t, \quad g_t \in \partial f(x_t, s_t).
\end{align}
It usually becomes unstable when the magnitude of $\alpha > 0$ increases.
It has been shown empirically that other methods are much less sensitive to large values of $\alpha$: in particular, this is known for the \emph{stochastic Polyak step size} \citep{Polyak1987,Asi2019,Loizou2021}: for a known lower bound $C_t \leq \inf_{z\in \R^d} f(z, s_t)$, let $g_t \in \partial f(x_t, s_t)$, then
\begin{align}
    \label{eqn:sps}\tag{\SPS{}}
    x_{t+1} = x_t - \min\{\alpha, \frac{f(x_t,s_t) - C_t}{\|g_t\|^2} \} g_t.
\end{align}
The above method allows for much larger $\alpha$ than \SGD{} across many problem instances in machine learning \citep{Schaipp2023,Schaipp2024a}.
%
Closely related to the Polyak step size is the non-negative Gauss-Newton step \citep{Orvieto2024}. In the stochastic setting, with $g_t \in \partial f(x_t, s_t)$, it is given by
\begin{align}
    \label{eqn:ngn}\tag{\NGN{}}
    x_{t+1} = x_t - \frac{\alpha}{1+\frac{\alpha}{2f(x_t,s_t)}\|g_t\|^2}g_t.
\end{align}

Finally, it has also been shown empirically that the stochastic proximal point method, given by
\begin{align}
    \label{eqn:spp}\tag{\SPP{}}
    x_{t+1} = \argmin_{y\in \R^d} f(y,s_t) + \frac{1}{2\alpha} \|y-x_t\|^2,
\end{align}
is less sensitive to the choice of $\alpha$ than \SGD{} \citep{Asi2019,Davis2019,Milzarek2024}.


In convex, nonsmooth optimization, we can usually derive bounds\footnote{For the concrete realization of such a bound, see \cref{thm:last-iterate}.} of the form
\begin{align*}
    \E[f(x_T) - f_\star] \leq \frac{\mathcal{T}_1}{\alpha} + \mathcal{T}_2(\alpha),
\end{align*}
where the term $\mathcal{T}_1/\alpha$ can be seen as a bias term; usually we have that $\mathcal{T}_1 \propto 1/T$.
However, the sensitivity to large step sizes is determined by the second term $\mathcal{T}_2(\alpha)$, reflecting the variance; for \SGD{} this term grows linearly with $\alpha$.

We will show that for other methods such as \SPS{}, \NGN{}, and \SPP{} the scaling  of $\mathcal{T}_2(\alpha)$ in $\alpha$ is much more benign. Further, we will show through experiments that the theoretical bound closely reflects the actual performance with respect to learning-rate tuning. This gives a theoretical explanation for the empirically observed improvements in stability for \SPS{}, \NGN{}, and \SPP{}.


\paragraph{Notation.}
%
Unless explicitly stated otherwise, $\|\cdot\|$ and $\iprod{\cdot}{\cdot}$ denote the standard Euclidean norm and its scalar product.
Throughout we will refer to $\alpha$ (and later $\alpha_t$) as the \emph{step size}, in the sense that this parameter has to be user-specified/tuned in practice. However, note that for the update rules of \eqref{eqn:sps} and \eqref{eqn:ngn} the \emph{effective} step size, i.e., the multiplier of $g_t$, is different from $\alpha$ in general. 
%%
\section{Related Work}\label{sec:related-work}
%
\paragraph{Model-based stochastic optimization.} The core concept of model-based optimization is, in each iteration, to construct a simplified model/surrogate of the objective, and to minimize this surrogate in a proximal fashion. Theory for this family of methods has been established in the seminal works of \citet{Asi2019} and \citet{Davis2019}. Concretely, \citet{Asi2019,Asi2019a} prove that models which share a certain lower-bound condition\footnote{For details, see condition (C.iv) in \citet{Asi2019}.} lead to methods that are stable, in the sense that the iterates are almost surely bounded for square-summable step sizes \citep[Cor.\ 3.5]{Asi2019}. However, this result does not quantify how performance degrades as the scale of the step sizes grows. Moreover, the convergence result that follows from stability \citep[Prop.\ 3.8]{Asi2019} is asymptotic. We refine this analysis by looking at non-asymptotic bounds, and their scaling with respect to the step size $\alpha$.
%%
\paragraph{The issue of learning-rate tuning.} Stochastic versions of the Polyak step size come in many different flavors \citep{Berrada2019,Loizou2021,Gower2022}. However, most of the practical implementations cap the effective step size with a hyperparameter $\alpha$ as depicted in equation~\eqref{eqn:sps}. This hyperparameter can be interpreted as a user-specified learning rate. Versions of \SPS{}, combinable with weight decay, momentum, and preconditioning, have shown promising empirical performance across many deep learning tasks, often being much less sensitive to the choice of $\alpha$ \citep{Schaipp2023,Schaipp2024a}. \SPS{} usually performs well across multiple orders of magnitude for $\alpha$. In this paper, we provide a theoretical analysis of this remarkable stability property of \SPS{}.\footnote{In \cref{sec:app:related-work}, we comment on why previous analyses such as the one by \citet{Loizou2021} do not fully reflect the empirical results.}

Similar to the Polyak step size, the non-negativity of loss functions in machine learning has been leveraged to derive Gauss-Newton type step sizes \citep{Orvieto2024,Islamov2025}. The resulting method \NGN{} has been proven to be similarly robust for deep learning tasks as Polyak-based methods \citep{Islamov2025}.

In this paper, we will focus mainly on \SPS{} and \NGN{}. However, we want to briefly summarize other efforts to design optimization algorithms that need less (or no) tuning of the step size (these are sometimes subsumed under the name ``parameter-free algorithms").
One line of work sets the learning rate on-the-fly by estimating a-priori unknown problem parameters based on the optimization trajectory \citep{Carmon2022,Defazio2023,Ivgi2023,Mishchenko2024}. Other parameter-free methods based on coin betting entirely remove the learning rate as hyperparameter \citep{Orabona2016,Orabona2017}. Recently, \citet{Kasimbeg2025} have conducted an empirical comparison of several of the methods mentioned above.

\paragraph{Stochastic proximal point.} For deterministic problems, it is well known that the proximal point method \citep{Martinet1970,Rockafellar1976a} converges for any sequence of step size, and potentially arbitrarily fast.\footnote{The caveat: a faster rate comes with a larger step size, which makes the subproblems in general harder to solve.}
In the stochastic case, \SPP{} has been analyzed in various settings \citep{Bertsekas2011,Boyd2014,Patrascu2017,Asi2019,Davis2019,Milzarek2024}. However, to the best of our knowledge, no clear theoretical advantage of \SPP{} has been established in the stochastic setting. Here, we prove an improvement in step-size stability for \SPP{}, and show how it is closely related to interpolation properties.
% so far: SPP is stable but SGD is not
% here: refined analysis, and relationship to interpolation

\begin{figure}[t]
    \begin{subfigure}{0.49\columnwidth}
       \centering
        \includegraphics[width=0.99\linewidth]{plots/models/logistic_scaling_obj.pdf}
    \end{subfigure}
    %
    \begin{subfigure}{0.49\columnwidth}
       \centering
        \includegraphics[width=0.99\linewidth]{plots/models/logistic_scaling_delta.pdf}
    \end{subfigure}
    \caption{Illustration of theory for $f(x, s_t) = \ln(1+\exp(-x)) + \max\{x-2, 0\}$ and $x_t=-3$. \textbf{(Left)} Next-iterate loss as a function of step size $\alpha$. \textbf{(Right)} Stability index $\delta_t$ as a function of $\alpha$. Vertical line marks best \SGD{} step size. For large $\alpha$, stable loss values coincide with benign ($\approx$ sub-linear) increase of $\delta_t$.}
    \label{fig:one}
\end{figure}

%%

\section{Theoretical Analysis}\label{sec:theory}
%%
We study model-based stochastic proximal point with step sizes $(\alpha_t)_{t\in \N} \subset \R_{> 0}$. Let $y\mapsto f_x(y,s)$ be a  model of $f(\cdot,s)$ around $x\in \R^d$, then the update is given by
\begin{align}\label{eqn:update}
    x_{t+1} = \argmin_{y \in \R^d} f_{x_t}(y,s_t) + \frac{1}{2\alpha_t}\|y-x_t\|^2.
\end{align}
%
We define $\delta_t$, called from now on the \textbf{stability index}, as
\begin{align}\label{eqn:def-delta}
\hspace{-2ex}\boxed{
    \delta_t := f(x_t,s_t) - f_{x_t}(x_{t+1},s_t) - \frac{1}{2\alpha_t}\|x_{t+1}-x_t\|^2.
} 
\end{align}
%
We will later show that $\delta_t$, more precisely its scaling with respect to $\alpha_t$, determines how stable the model-based update~\eqref{eqn:update} is to the choice of the step size $\alpha_t$. \cref{fig:one} illustrates this phenomenon on a toy example. Before we derive  $\delta_t$ for concrete choices of $f_x(\cdot,s)$ (see \cref{sec:models}), we start by presenting the role of $\delta_t$ in suboptimality (resp.\ stationarity) bounds for the convex (resp.\ weakly convex) case. 

\subsection{Convex Case}\label{sec:theory-convex}
%
In this sub-section, we make the following model assumptions: for every $s\in \mathcal{S}$, it holds:
\begin{enumerate}[label=(A\arabic*)]
    \item\label{item:A1} The mapping $y\mapsto f_x(y,s)$ is convex.
    \item\label{item:A2} We have $f_x(x,s) = f(x,s)$ for all $x\in \R^d$ and $f_x(y,s) \leq f(y,s)$ for all $y\in \R^d$.
\end{enumerate}
%
We will show later that if $f(\cdot,s)$ is convex, for many methods considered here, the above two assumptions are satisfied (see \cref{sec:models}).
%
Further, from the definitions \eqref{eqn:def-delta} and \eqref{eqn:update}, it follows immediately that the stability index is always non-negative under \ref{item:A2}.
\begin{lemma}\label[lemma]{lem:delta-nonneg}
    If the model satisfies $f_x(x,s) = f(x,s)$, in particular if \ref{item:A2} holds, then $\delta_t \geq 0$ holds for any $x_t \in \R^d,s_t \in \mathcal{S}$ and any $\alpha_t>0$.
\end{lemma}

Note that the model-based framework implicitly anchors the step size of all methods onto the same scale, therefore making them comparable: this is due to update \eqref{eqn:update}, and the first part of \ref{item:A2}.


Next, we state the base lemma for our theoretical analysis.
\begin{lemma}\label[lemma]{lem:base}
Let $1\leq k\leq T$ and let $u\in \R^d$ be measurable with respect to $x_k$. Let the iterates $(x_t)_{t\in \N}$ be generated by \eqref{eqn:update}. Assume that \ref{item:A1}-\ref{item:A2} hold.
Denoting $D_t := \|x_t-u\|^2$, we have
    \begin{align*}
        \sum_{t=k}^T \alpha_t \E[f(x_t) - f(u)] \leq
        %\frac12\E\|x_k - u\|^2 - \frac12\E\|x_{T+1} - u\|^2
        \frac12\E[D_k-D_{T+1}]
        + \sum_{t=k}^T \alpha_t \E[\delta_t].
    \end{align*}
\end{lemma}
%
\begin{proof}
    Under \ref{item:A1}, \citet[Lem.\ 3.7]{Asi2019} show
    \begin{align*}
        \frac12 \|x_{t+1} - u\|^2 \leq \frac12 \|x_t-u\|^2   -
        \frac12 \|x_{t+1}-x_t\|^2 \\
        - \alpha_t[\textcolor{OtherBlue}{f_{x_t}(x_{t+1},s_t) - f_{x_t}(u,s_t)}].
    \end{align*}
    Now, we decompose
    \begin{align*}
        &\hspace{-2ex} \textcolor{OtherBlue}{f_{x_t}(x_{t+1},s_t) - f_{x_t}(u,s_t)} =  \\
        &\underbracket{f(x_t,s_t) - f_{x_t}(u,s_t)}_{=:~U_1} + \underbracket{f_{x_t}(x_{t+1},s_t) - f(x_t,s_t)}_{=:~U_2}.
    \end{align*}
    Due to \ref{item:A2} we have $U_1 \geq f(x_t,s_t) - f(u,s_t)$.
    For the other term, use \eqref{eqn:def-delta} to obtain
    \begin{align*}
        -\alpha_t U_2 - \frac12\|x_{t+1}-x_t\|^2 = \alpha_t \delta_t. 
    \end{align*}
    Altogether, we get
    \begin{align}\label{eqn:single-step-base-ineq}
        %\frac12 \|x_{t+1} - u\|^2 \leq \frac12 \|x_t-u\|^2  - \alpha_t[f(x_{t},s_t) - f(u,s_t)] + \alpha_t \delta_t.
        \frac12 (D_{t+1}-D_t) \leq - \alpha_t[f(x_{t},s_t) - f(u,s_t)] + \alpha_t \delta_t.
    \end{align}
    Then apply conditional expectation, and sum over $t=k,\ldots,T$.
\end{proof}
%
\begin{theorem}[Average-iterate bound]\label{thm:avg-iterate}
    Assume that \ref{item:A1}-\ref{item:A2} hold. Let $T\in \N$ and let the iterates $(x_t)_{t\in \N}$ be generated by \eqref{eqn:update}. Let $x_\star \in \R^d$ and $D:=\|x_1-x_\star\|$.
    Then, it holds, for both $\bar f_T = \E[f((\sum_{t=1}^T \alpha_t)^{-1} \sum_{t=1}^T \alpha_t x_t)]$ and $\bar f_T = \min_{t=1,\ldots,T} \E[f(x_t)]$ that
    \begin{align*}
        \bar f_T - f(x_\star) \leq \frac{D^2}{2\sum_{t=1}^T \alpha_t} + \frac{\sum_{t=1}^T \alpha_t \E[\delta_t]}{\sum_{t=1}^T \alpha_t}.
    \end{align*}
\end{theorem}
%
\begin{proof}
    Apply \cref{lem:base} with $k\to 1$ and $u\to x_\star$. Then, divide by $\sum_{t=1}^T \alpha_t$ (and apply Jensen's inequality).
\end{proof}
%
\begin{theorem}[Last-iterate bound]\label{thm:last-iterate}
    Assume that \ref{item:A1}-\ref{item:A2} hold. Let $T\in \N$ and let the iterates $(x_t)_{t\in \N}$ be generated by \eqref{eqn:update}. Let $x_\star \in \R^d$ and $D:=\|x_1-x_\star\|$. Then, it holds
    \begin{align*}
        &\E[f(x_T) - f(x_\star)] \leq \frac{D^2}{2\sum_{t=1}^T \alpha_t} + \frac{ \sum_{t=1}^T \alpha_t\E[\delta_t]}{\sum_{t=1}^T \alpha_t} \\
        &\hspace{10ex} + \sum_{k=1}^{T-1} \frac{\alpha_k}{\sum_{t=k+1}^{T}\alpha_t} \Big(\frac{1}{\sum_{t=k}^{T}\alpha_t} \sum_{t=k}^T \alpha_t \E[\delta_t]\Big).
    \end{align*}
\end{theorem}
%
\begin{proof}
    Follows identically to the proof of \citet[Thm.\ 10]{Defazio2023a} (or \citet[Thm.\ 3.1]{Schaipp2025}), using \cref{lem:base} instead of \citet[Lem.\ E.1]{Schaipp2025}, and replacing $\frac12 \eta_t^2 \E\|g_t\|^2$ by $\alpha_t \E[\delta_t]$.
\end{proof}
%
Throughout, we will denote the right hand side of the bound in \cref{thm:avg-iterate,thm:last-iterate} as follows:
\begin{align*}
    \bar f_T - f(x_\star) \leq \Omega_T^{\text{avg}}, \quad \E[f(x_T) - f(x_\star)] \leq \Omega_T^{\text{last}}.
\end{align*}
%
\paragraph{Illustration of \cref{thm:last-iterate}.} Before deriving the value of~$\delta_t$ for the methods of interest, we illustrate how the bound $   \Omega_T^{\text{last}}$ in \cref{thm:last-iterate} (analogously \cref{thm:avg-iterate}) reflects the aspect of stability with respect to large step sizes $\alpha_t$ (see \cref{fig:illustration-nu}). For this, we fix a constant step size $\alpha_t = \alpha$. 

As discussed above, the scaling behavior in $\alpha$ depends crucially on the scaling of $\E[\delta_t]$ in $\alpha$. Purely for the sake of illustration, we now make the simplification $\Delta_t:=\E[\delta_t] = \alpha^{\nu}$ for some $\nu \geq 0$.
Denote $H_{T} := \sum_{s=1}^{T} \frac{1}{s}$. Then, with $\alpha_t = \textcolor{OtherBlue}{\alpha} $ and $\Delta_t= \alpha^{\nu}$, we have
\begin{align*}
    \Omega_T^{\text{last}} = \frac{D^2}{2\textcolor{OtherBlue}{\alpha}T} + {\textcolor{OtherBlue}{\alpha}}^\nu [1+H_{T-1}].
\end{align*}
%

\takeaway{The stability of the bound 
$\Omega_T$ 
with respect to large learning rates $\alpha$ crucially depends on the scaling of $\Delta_t = \E[\delta_t]$ (as illustrated with our simplification $\Delta_t= \alpha^{\nu}$, see \cref{fig:illustration-nu}).}

\begin{figure}
    \begin{subfigure}{0.49\columnwidth}
        \centering
        \includegraphics[width=0.95\linewidth]{plots/bound/sweep_nu_scaling.pdf}
        \caption{Illustration of \cref{thm:last-iterate}}
        \label{fig:illustration-nu}
    \end{subfigure}
    %
    \begin{subfigure}{0.49\columnwidth}
    \centering
    \includegraphics[width=0.95\linewidth]{plots/delta_spp/pepit_bound.pdf}
    \caption{\texttt{PEPit} verification}
    \label{fig:spp-pepit-delta}
\end{subfigure}
\caption{\textbf{(Left)} Stability with respect to learning-rate $\alpha$ under different scalings of $\Delta_t$. \textbf{(Right)} Values of $\Lambda(\alpha)$ \textbf{(blue)} computed with \texttt{PEPit}, for different Lipschitz constants $G$. We plot the upper bound $f_s(x) -\inf f_s$ for the worst-case instance in \textbf{(green)}, and the corresponding $\delta^{\SGD{}} = \frac12 \alpha G^2$ \textbf{(grey)}.}
\end{figure}
%%
%
%
\subsection{Weakly Convex Case}\label{sec:theory-weakly-convex}
%
This section shows that the role of $\E[\delta_t]$ is not limited to the convex case.
For this, we first define a function $\phi:\R^d\to \R$ to be $\rho$-weakly convex for $\rho \geq 0$ if $\phi + \frac{\rho}{2}\|\cdot\|^2$ is convex. In this case, for any subgradient $u \in \partial \phi(x)$ it holds $\phi(y) \geq \phi(x) + \iprod{u}{y-x} -\frac{\rho}{2}\|y-x\|^2$ for all $y\in \R^d$ \citep[Lem.\ 2.1]{Davis2019}.
If $\phi$ is $\rho$-weakly convex, its \emph{Moreau envelope} \citep{Moreau1965} is given by
\begin{align*}
    \env{\phi}^\alpha (x) := \min_{y\in \R^d} \phi(y) + \frac{1}{2\alpha}\|y-x\|^2, \quad \alpha < 1/\rho.
\end{align*}
%
The Moreau envelope is differentiable, and $\|\nabla \env{\phi}^\alpha (x)\|$ can be interpreted as a measure of near-stationarity of $\phi$ at $x$ (see \citet[Section 2.2]{Davis2019}).

We replace \ref{item:A1} and \ref{item:A2} with their natural generalization:
\begin{enumerate}[label=(\~{A}\arabic*)]
    \item\label{item:tA1} For any $x\in \R^d$ and $s\in \mathcal{S}$, the mapping $y\mapsto f_x(y,s)$ is $\lambda$-weakly convex for some $\lambda\geq 0$.
    \item\label{item:tA2} We have $f_x(x,s) = f(x,s)$ for all $x\in \R^d$. Further, for all $s\in \mathcal{S}$ there exists $\rho_s \geq 0$ such that for all $x,y\in \R^d$ we have $f_x(y,s) \leq f(y,s) + \frac{\rho_s}{2}\|x-y\|^2$. Assume that $\rho := \E_s[\rho_s] < +\infty $.
\end{enumerate}
%
Note that if \ref{item:tA1} and \ref{item:tA2} are satisfied with $\lambda=\rho_s=0$, then \ref{item:A1} and \ref{item:A2} are satisfied.
Further, \ref{item:tA1}-\ref{item:tA2} imply that $f=\E[f(\cdot,s)]$ itself is $(\rho+\lambda)$-weakly convex \citep[Lem.\ 4.1]{Davis2019}.


Below, we state analogous results to \cref{lem:base,thm:avg-iterate} for the weakly convex case. Proofs are deferred to \cref{sec:app:proofs-weakly-convex} in the Appendix.

\begin{lemma}\label[lemma]{lem:weakly-convex-base}
    Assume that \ref{item:tA1} and \ref{item:tA2} hold. Let $\alpha_t \in (0, 1/\lambda)$. Let $x_t\in \R^d$ and denote by $\E_t$ the conditional expectation with respect to the filtration generated by $\{x_1,\ldots,x_t\}$. Let $\bar{\rho} > \rho + \lambda$ and let $\hat{x}_t = \prox{(1/\bar{\rho})f}(x_t)$. Then, we have
    \begin{align}\label{eqn:base-ineq-weakly-convex}
    \begin{split}
        &\hspace{-4ex}\E_t\|x_{t+1} - \hat{x}_t\|^2 \leq \|x_{t} - \hat{x}_t\|^2 \\
        &\hspace{0ex}- \frac{\alpha_t(\bar{\rho}-\rho-\lambda)}{1-\alpha_t\lambda}\|x_{t} - \hat{x}_t\|^2
        + \frac{2\alpha_t}{1-\alpha_t\lambda} \E_t[\delta_t].
    \end{split}
    \end{align}
\end{lemma}
%
\begin{theorem}\label{thm:weakly-convex}
    Assume that \ref{item:tA1} and \ref{item:tA2} hold, which in particular implies that $f$ is $(\rho+\lambda)$-weakly convex. Let $1\leq k \leq T$ and let the iterates $(x_t)_{t\in \N}$ be generated by \eqref{eqn:update}. Let $\alpha_t \in (0, 1/\bar{\rho})$ for all $t\in \N$ for some $\bar{\rho} > \rho + \lambda$. Denoting $\mathcal{E}_t:= \E[\env{f}^{1/\bar{\rho}}(x_t)]$, it holds
    \begin{align*}
        &\sum_{t=k}^{T} \frac{\alpha_t}{1-\alpha_t\lambda} \E\|\nabla \env{f}^{1/\bar{\rho}}(x_t)\|^2 \leq \frac{2\bar{\rho}}{\bar{\rho}-\rho-\lambda} \big[\mathcal{E}_k - \mathcal{E}_{T+1}\big] \\
        &\hspace{5ex}+ \frac{2\bar{\rho}^2}{\bar{\rho}-\rho-\lambda}\sum_{t=k}^{T}\big(\frac{\alpha_t}{1-\alpha_t \lambda} \E[\delta_t]\big).
    \end{align*}
\end{theorem}
%
The above bound can be easily turned into a convergence rate for the near-stationarity measure $\E\|\nabla \env{f}^{1/\bar{\rho}}(\cdot)\|^2$, for example see \citep[Thm.\ 4.3]{Davis2019}. We can see that the role of $\E[\delta_t]$ in \cref{thm:weakly-convex} is analogous to the  previous results in \cref{sec:theory-convex}.
%
\section{Model Choices}\label{sec:models}
%
In this section, we derive the stability index \eqref{eqn:def-delta} for classical model choices.
Proofs are deferred to \cref{sec:app:proofs-deltas} in the Appendix. Let the \emph{regularized model} be given by
\begin{align}\label{eqn:psi-delta}
\Psi_t:=f_{x_t}(x_{t+1},s_t) + \frac{1}{2\alpha_t}\|x_{t+1}-x_t\|^2,
    %~ \text{hence}~ \Psi_t= f(x_t,s_t) - \delta_t.
\end{align}
hence $\Psi_t= f(x_t,s_t) - \delta_t$.
%
We assume from now on that $f(\cdot,s)$ is $\eta_s$-weakly convex, that is, there exists $\eta_s \geq 0$ such that for all $x,y \in \R^d$ it holds
\begin{align}\label{eqn:weak-convexity}
    f(y,s) - f(x,s) \geq \iprod{g}{y-x} - \frac{\eta_s}{2} \|y-x\|^2
\end{align}
for all $g\in \partial f(x,s)$ and all $s\in \mathcal{S}$.
%Convexity is sufficient to guarantee \ref{item:A1} and \ref{item:A2} for \SGD{} and \SPS{} \citep{Asi2019}.

%Denote the expectation conditioned on the filtration generated by the iterates $\{x_1,\ldots,x_t\}$ with $\E_t[\cdot]$. 
%
\subsection{Linear Model (\SGD{})} \label{sec:linearmodel}
%
If $f_x(y,s)=f(x,s) + \iprod{g}{y-x}$ for some $g\in \partial f(x,s)$, then $x_{t+1} = x_t -\alpha_t g_t$ and using this model in~\eqref{eqn:psi-delta} gives
$\Psi_t = f(x_t, s_t) - \frac{\alpha_t}{2}\|g_t\|^2.$
Therefore, we have $\delta_t^{\SGD{}} = \frac{\alpha_t}{2}\|g_t\|^2$ which recovers \citet[Thm.\ 3.1]{Schaipp2025}.
It is easy to see that \ref{item:tA1} holds with $\lambda=0$ and \ref{item:tA2} holds with $\rho_s=\eta_s$ due to~\eqref{eqn:weak-convexity}.
%%%%%%%
%
\subsection{Truncated Model (\SPS{})}\label{sec:spsmodel}
%
Given access to a lower bound on our stochastic function $C_s\leq \inf_z f(z,s)$ we can define a \emph{truncated model}
\begin{equation}\label{eq:truncated-model}
    f_x(y,s)=\max\{f(x,s) + \iprod{g}{y-x}, C_s\}
\end{equation}  
for some $g\in \partial f(x,s)$.
%This model $f_x(\cdot,s)$ is piecewise linear, hence convex and \ref{item:A1} holds.~\ref{item:A2} holds due to convexity \eqref{eqn:convexity} and the fact that $C_s \leq f(y,s)$ for all $y\in \R^d$.
This model satisfies \ref{item:tA1} with $\lambda=0$ and \ref{item:tA2} with $\rho_s=\eta_s$ \citep[Lem.\ 2]{Schaipp2023}. 
This truncated model also gives a stability index that is always smaller or equal than the stability index of \SGD{}.

%\begin{restatable}{lemma}{lemdeltasps}\label{lem:delta-sps}
\begin{lemma}\label[lemma]{lem:delta-sps}
    Let  $s \mapsto C_s$ be a mapping such that $C_s \leq \inf_z f(z,s).$ Using the truncated model~\eqref{eq:truncated-model} in~\eqref{eqn:update} results in the \SPS{} (Stochastic Polyak Stepsize) method
    \begin{align}\label{eqn:def-tau}
        x_{t+1} = x_t - \tau_t g_t, \quad  \tau_t:= \min\{\alpha_t, \frac{f(x_t,s_t)-C_{t}}{\|g_t\|^2}\}.
    \end{align}
    Its stability index is $\delta^{\SPS{}}_t = \tau_t[1- \frac{\tau_t}{2\alpha_t}]\|g_t\|^2$. In particular, $\delta^{\SPS{}}_t \; \leq\; \delta^{\SGD{}}_t$ and
    \begin{align} \label{eq:deltaSPS-upper}
       \boxed{ \delta^{\SPS{}}_t \leq  \min\{\alpha_t\|g_t\|^2, f(x_t,s_t)-C_t \}. } 
    \end{align}
\end{lemma}
%\end{restatable}


% \takeaway{$\delta_t^{\SPS{}}$ scales sub-linearly with $\alpha_t$, as it is capped by $f(x_t,s_t)-C_t$. We further have that $\delta_t^{\SPS{}} \leq \delta_t^{\SGD{}}$ holds for \emph{any} choice of $C_t$.
% }
We can also bound the expected stability index, which is the final quantity that appears in our convergence theorems.
Using \cref{lem:exp-of-min} together with~\eqref{eq:deltaSPS-upper}, we have 
\begin{align*}
    \E[\delta_t^{\SPS{}}] \leq \min\{\alpha_t \E\|g_t\|^2, \E[f(x_t)- C_{t}]\}.
\end{align*}
The second term in the $\min$-operator can be decomposed as
\begin{align*}
    &\hspace{-4ex}\E[f(x_t)- C_{t}] = \E[f(x_t) - f(x_\star)] \\ +
    &\underbracket{f(x_\star)-\E_s[\inf_z f(z,s)]}_{=:~\sigma^2} + 
    \underbracket{\E_s[\inf_z f(z,s) - C_s]}_{=:~\epsilon_{\text{lb}}}.
\end{align*}
%
Here we call $\sigma^2$ the \emph{interpolation constant} and $\epsilon_{\text{lb}}$ the \emph{lower-bound estimation error}.
%
\takeaway{We always have $\delta_t^{\SPS{}} \leq \delta_t^{\SGD{}}$ (assuming both methods are at the same $x_t$ and $s_t$). The gain in stability for \SPS{} is determined by the interpolation constant, and the lower-bound estimation error.}
%
\subsection{Square-root model (\NGN{})} \label{sec:ngnmodel}
For non-negative loss functions, we can use the \NGN{} (Non-negative Gauss Newton) method proposed by \citet{Orvieto2024}.
They propose to rewrite $f(x) = \E[f(x,s)]=\E[\sqrt{f(x,s)}^2]$ and, for given $s\in \mathcal{S}$, to linearize $y \mapsto \sqrt{f(y,s)}$ around $x$, followed by the square operation. 
%
This leads us to the following model: for $g \in \partial f(x,s)$, define
% \footnote{
% To define this model, we act as if the subdifferential of $f(\cdot,s)$ allows for the chain rule.
% }
\begin{align}\label{eqn:ngn-model}
    f_{x}(y,s) = \big(\sqrt{f(x,s)} + \frac{1}{2\sqrt{f(x,s)}}\iprod{g}{y-x} \big)^2.
\end{align}
This \NGN{} model also enjoys a favorable stability index.
%\begin{restatable}{lemma}{lemdeltangn}\label{lem:delta-ngn}
\begin{lemma}\label[lemma]{lem:delta-ngn}
    Let $f(x,s) \geq 0$ for all $x\in \R^d,~s\in \mathcal{S}$. Using model~\eqref{eqn:ngn-model} results in the \NGN{} method \citep{Orvieto2024}
    \begin{align}\label{eq:NGNstep}
    x_{t+1} = x_t - \gamma_t g_t, \quad \gamma_t := \frac{\alpha_t}{1+\frac{\alpha_t}{2f_t}\|g_t\|^2}.
    \end{align}
Its stability index is given by
\begin{equation*}
    \boxed{
    \delta_t^{\NGN{}}\; =\; \frac{\gamma_t}{2} \|g_t\|^2 \; \leq \; \delta_t^{\SGD{}}.
    }
\end{equation*}
\end{lemma}
%\end{restatable}

Not only do we have $\delta_t^{\NGN{}} \leq \delta_t^{\SGD{}}$, but further if $\alpha_t\to + \infty$ we have $\gamma_t \to 2\frac{f(x_t,s_t)}{\|g_t\|^2}$. Consequently, $\delta_t^{\NGN{}}$ does not scale linearly with $\alpha_t$, which suggests that the method should be very stable for large step sizes.
Also note the relation to the Polyak step size when $C_t=0$ (up to a factor of two). 

\takeaway{It always holds $\delta_t^{\NGN{}} \leq \delta_t^{\SGD{}}$ and further $\delta_t^{\NGN{}}$ scales sub-linearly with respect to $\alpha_t \to +\infty$.}

%%%%%%%%%%%%%%%%%%%%%%%%%%%%%%%%%%%%%%%%%%%%%%%%%%%%
%%% A1 and A1 assuming that sqrt(f(x,s)) is convex
%%%%%%%%%%%%%%%%%%%%%%%%%%%%%%%%%%%%%%%%%%%%%%%%%%%%
% It remains to check when will the model Assumptions (A1)   and (A2) hold for the \NGN{} model.
% \begin{lemma}
%    Consider the model~\eqref{eqn:ngn-model}. Assumption (A1) holds the mapping $y\mapsto f_x(y,s)$ is convex as a composition of convex and linear mappings. Assumption (A2) holds iff $f(x,s) \geq 0$ and $h(x) : = \sqrt{f(x,s)}$ satisfies
%    \begin{equation} \label{eq:symmetricconvex}
%         h(y) \geq |h(x) + \iprod{s(x)}{ y-x} |,
%    \end{equation}
%    where $s(x) \in \partial h(x).$ 
% \end{lemma}
% \begin{proof}
% Regarding \ref{item:A2}, we clearly have $f_x(x,s) = f(x,s)$. 
% Furthermore, since $f(x,s)$ is non-negative and $\sqrt{f(x,s)}$ is convex, by  Lemma~\ref{lem:squareconv} we have that $f(x,s) = (\sqrt{f(x,s)})^2$ is convex, and thus $f(x,s)$ has a subgradient and our model~\eqref{eqn:ngn-model} is well defined.

% For the second condition in Assumption (A2), since both sides of~\eqref{eq:symmetricconvex} are positive, taking the square of both sides preserves the inequality giving
% \begin{align}
%     f(y,s) = h(y)^2  & \geq (h(x) + \iprod{s(x)}{ y-x})^2 \nonumber \\
%     & = (\sqrt{f(x,s)} + \iprod{s(x)}{ y-x})^2 \label{eq:aewdadwdadA}
% \end{align}
% By the chain rule of subgradients by \citet[Corollary 16.72]{bauschke2011convex} we have 
% that $s(x) := g(x)/2f(x,s)$ 
% is a subgradient of $\sqrt{f(x,s)}$, where $g(x)$ is a subgradient of $f(x,s)$. Consequently using this subgradient $s(x) := g(x)/2f(x,s)$  in~\eqref{eq:aewdadwdadA} gives
% \begin{align}
%     f(y,s)   & \geq  (\sqrt{f(x,s)} + \frac{1}{2f(x,s)}\iprod{g(x)}{ y-x})^2 \label{eq:aewdadwdadAsdss}
% \end{align}
% which is exactly the second condition in Assumption (A2).
% \end{proof}
% This \emph{symmetric} convexity condition in~\eqref{eq:symmetricconvex} holds for instance for $h(x) = \| Ax-b\|$ for any norm $\|\cdot\|.$
%%%%%%%%%%%%%%%%%%%%%%%%%%%%%%%%%%%%%%%%%%%%%%%%%%%%
%%% END
%%%%%%%%%%%%%%%%%%%%%%%%%%%%%%%%%%%%%%%%%%%%%%%%%%%%

\paragraph{On Assumptions \ref{item:tA1} and \ref{item:tA2}.} \ref{item:tA1} holds with $\lambda=0$ due to the convexity of a composition of convex and linear mappings. Regarding \ref{item:tA2}, we clearly have $f_x(x,s) = f(x,s)$. However, the second part of \ref{item:tA2} is less trivial. It is equivalent to the condition that, for some $\rho_s\geq 0$, it holds
\begin{align*}
    f(y,s) + \tfrac{\rho_s}{2}\|y-x\|^2 \geq f(x,s) + \iprod{g}{y-x} + \frac{\iprod{y-x}{g}^2}{4f(x,s)} 
\end{align*}
for all $x,y\in \R^d$, all $s\in \mathcal{S}$ and all $g\in \partial f(x,s)$.
Assuming a global upper bound on $\|g\|$ and lower bound on $f(x,s)$ this can be achieved with $\rho_s > \eta_s$. For the convex case, where we require $\rho_s=0$, the above condition can be satisfied at least on a bounded domain, if $f(\cdot,s)$ is $\alpha$-exp-concave (see \citet[Def.\ 4.1]{Hazan2019a}. This follows from \citet[Lem.\ 4.3]{Hazan2019a}, considering that $\frac{1}{2f(x,s)} \leq \frac{1}{2\inf_z f(z,s)}$ and the latter term can be made arbitrarily small by adding a constant to $f(\cdot, s)$. 

\paragraph{Extensions of \NGN{}.} We show in \cref{sec:app:ngn-generalizations} of the Appendix how to extend the key idea of \NGN{} to general composition of maps. In particular, we show that under mild assumptions on these maps, it always holds that the resulting stability index is smaller than the one of \SGD{}.

%
%
\subsection{Exact Model (\SPP{})} \label{sec:sppmodel}
%
Using the exact model, that is $f_x(y,s) = f(y,s)$, recovers the \SPP{} method. Assume for this part that $\bar{\eta} := \sup_{s\in \mathcal{S}} \eta_s < +\infty$. Then, $f(\cdot,s)$ is $\bar{\eta}$-weakly convex for all $s\in \mathcal{S}$, and update \eqref{eqn:update} is well-defined for $\alpha_t \in (0, 1/\bar{\eta})$. In particular, \ref{item:tA1} is satisfied with $\lambda=\bar{\eta}$, and \ref{item:tA2} with $\rho_s=0$.

%\begin{restatable}{lemma}{lemdeltaspp}\label{lem:delta-spp}
\begin{lemma}\label[lemma]{lem:delta-spp}
    Assume that $f(\cdot,s)$ is $\bar{\eta}$-weakly convex and $\alpha_t \in (0, 1/\bar{\eta})$. Then, the stability index for the exact model $f_x(y,s) = f(y,s)$ is bounded by
    \begin{align*}
    \boxed{
        \delta_t^{\SPP{}} \leq \min\{\tfrac{\alpha_t}{2(1-\alpha_t \bar{\eta})}\|g_t\|^2, f(x_t,s_t) - \inf_{y\in \R^d} f(y,s_t)\}.}
    \end{align*}
    In particular, if $\bar{\eta} = 0$, then $\delta_t^{\SPP{}} \leq \delta_t^{\SGD{}}$ for all $\alpha_t > 0$.
\end{lemma}
%\end{restatable}

%
Using \cref{lem:exp-of-min} reveals the relationship to the interpolation constant $\sigma^2$: 
\begin{align*}
    \E[\delta_t^{\SPP{}}] \leq \min\{\tfrac{\alpha_t}{2(1-\alpha_t \bar{\eta})}\E\|g_t\|^2, \E[f(x_t)-f(x_\star)] + \sigma^2\}.
\end{align*}


%
%% Original: 
% \takeaway{$\delta_t^{\SPP{}}$ scales sub-linearly with $\alpha_t$. Much like \SPS{}, the bound on $\E[\delta_t^{\SPP{}}]$ for large $\alpha_t$ decreases with the interpolation constant $\sigma^2$.
% }
\takeaway{In the convex case, $\delta_t^{\SPP{}}$ scales sub-linearly with $\alpha_t$. The bound on $\E[\delta_t^{\SPP{}}]$ is also capped by sub-optimality and  the interpolation constant $\sigma^2$.
}

\begin{remark}
    Note that for finite-sum problems, where $f(\cdot,s)$ defines a mini-batch loss, the choice of the batch size influences the value of $\sigma^2$ (for example, consider linear regression, where $\inf_z f(z,s)=0$ as long as the batch size is smaller than the dimension $d$, but $\inf_z f(z,s)$ can be non-zero for large batch sizes).
    The above shows that the batch size has an impact on the stability of \SPP{}; indeed we observe this in the experiments below, and is further confirmed by the empirical results of \citet{Milzarek2024}.
\end{remark}

%
\paragraph{Verification with PEP.} In contrast to the other methods, for \SPP{} we can only derive an \emph{upper bound} of $\delta_t$. In this paragraph, we verify whether our bound for $\delta_t^{\SPP{}}$ is tight in a worst-case sense (for the convex case, $\bar{\eta}=0$). For this, fix $x \in \R^d$ and let $x_+ = \argmin_{y\in \R^d} f(y,s) + \frac{1}{2\alpha}\|y-x\|^2$ be the \SPP{} update for fixed $t\in \N$.
Denote by $f_s$ the mapping $y \mapsto f(y,s)$.
%Let $\env{f_s}^\alpha(x)$ be the Moreau envelope, defined by $\env{f_s}^\alpha(x) := \min_{z \in \R^d} f_s(z) + \frac{1}{2\alpha}\|z-x\|^2$ (see \citet{Beck2017} for more details). 
Let $\delta(\alpha) = f_s(x) - \env{f_s}^\alpha(x)$. Denote by $\mathcal{C}$ the class of lower-bounded convex functions. Assume that $u_\star \in \argmin_z f_s(z)$. For $G> 0$, and $g\in \partial f_s(x)$, we are interested in the quantity
\begin{align*}
    \Lambda(\alpha) := \max_{f_s \in \mathcal{C}} \delta(\alpha) ~ \text{s.t.\ }~ f_s(x) - \min f_s \leq 1,~ \|g\| \leq G.
\end{align*}
It is easy to see that $\env{f_s}^\alpha(x) \searrow \inf f_s$ with $\alpha\to \infty$, and hence $\Lambda(\alpha)\leq f_s(x) - \inf f_s$. However, this bound is not necessarily tight, especially for small $\alpha$.

Therefore, we compute $\Lambda(\alpha)$ using the PEP (Performance Estimation Problem) framework \citep{Drori2014,Taylor2017} implemented in \texttt{PEPit} \citep{Goujaud2024}.
\cref{fig:spp-pepit-delta} reveals that for the worst-case example the bound
%\begin{align*}
    $\Lambda(\alpha) \leq \min\{\frac{\alpha}{2} G^2, f_s(x) - \inf f_s(x)\}$
%\end{align*}
is indeed very tight across all $\alpha > 0$; only in the region where the transition of the minimum happens, we can observe that $\Lambda(\alpha)$ is slightly smaller.
%
\begin{remark}[Practicality of \SPP{}]
    \SPP{} is often considered a theoretical method only, as the update in general cannot be solved in closed form. However, \citet{Milzarek2024} show how the update subproblem can be solved in competitive runtime for specific problem structures that naturally appear in statistical learning (generalized linear models). Their technique allows for weakly convex loss functions, mini-batching as well as nonsmooth regularization functions.
\end{remark}
%

\subsection{Summary}
%
We summarize the main aspects of our theoretical analysis:

\hspace{2ex}\textbf{1)} For the question of training stability with respect to large learning rates, the key quantity is $\delta_t$ and how it scales (in expectation) with respect to the step size $\alpha_t$.

\hspace{2ex}\textbf{2)} We derive the value of $\delta_t$ for four methods, \SGD{}, \SPS{}, \NGN{}, and \SPP{}. We show that \SPS{}, \NGN{}, and \SPP{} have stability index $\delta_t$ that is smaller or equal to that of \SGD{}.  Indeed, for \SGD{} the scaling of $\delta_t$ wrt.\ $\alpha_t$ is linear, whereas for \NGN{}, \SPS{} and \SPP{} this linear growth is flattened. This is illustrated in \cref{fig:toy-delta-illustration}.  As a by-product of our analysis of $\delta_t$, by plugging our bounds for  \NGN{} and \SPS{} into~\cref{thm:avg-iterate,thm:last-iterate}, we get the first convergence analysis in the convex and non-smooth setting for these methods (see~\cref{sec:app:conv-comparison}). 

\hspace{2ex}\textbf{3)} We stress that (weak) convexity of $f(\cdot,s)$ is \emph{not} required to derive the above expressions of $\delta_t$ (with the exception of \SPP{}). Convexity is only required for the argument how $\delta_t$ impacts the suboptimality bound \cref{thm:avg-iterate,thm:last-iterate}. Our extension to weakly convex problems shows that $\delta_t$ is informative of stability also for non-convex problems -- and indeed our experiments below confirm that.

We state again the key quantity $\delta_t$ for each of the methods we considered.\footnote{For \SPP{} we state the formula for the convex case, $\bar{\eta}=0$.} While we focus on these four models, our analysis can be extended to other model-based methods.
\begin{align}\label{eqn:delta-summary}
    \begin{split}
        \delta_t^{\SGD{}} &= \tfrac{\alpha_t}{2} \|g_t\|^2, \\
        \delta_t^{\SPS{}} &= \tau_t[1-\frac{\tau_t}{2\alpha_t}]\|g_t\|^2, \quad
        % \big(\tau_t = \min\{\alpha_t, \tfrac{f(x_t,s_t)-C_{t}}{\|g_t\|^2}\}\big)
        (\tau_t \text{ from } \eqref{eqn:def-tau})
        ,\\
        \delta_t^{\NGN{}} &= \tfrac{\gamma_t}{2} \|g_t\|^2, \quad
        % \big(\gamma_t = \frac{\alpha_t}{1+\frac{\alpha_t}{2f(x_t,s_t)}\|g_t\|^2}\}\big),
        (\gamma_t \text{ from } \eqref{eq:NGNstep}),\\
        \delta_t^{\SPP{}} &\leq \min\big\{\tfrac{\alpha_t}{2}\|g_t\|^2, f(x_t,s_t) - \inf_{y\in \R^d} f(y, s_t) \big\}.
    \end{split}
\end{align}
%
The above summary on how $\delta_t$ scales with $\alpha_t$ for different methods can be seen even clearer for the special case of quadratic functions where the above expressions simplify  further (see \cref{sec:app:quadratics})
%
\begin{figure}[t]
    \begin{subfigure}{0.49\columnwidth}
       \centering
        \includegraphics[width=0.99\linewidth]{plots/models/logistic_delta_sgd.pdf}
        \caption{\SGD{}}
    \end{subfigure}
    %
    \begin{subfigure}{0.49\columnwidth}
       \centering
        \includegraphics[width=0.99\linewidth]{plots/models/logistic_delta_ngn.pdf}
        \caption{\NGN{}}
    \end{subfigure}
    \caption{Illustration of the stability index $\delta_t$ for \SGD{} (left) and \NGN{} (right), for the function from \cref{fig:one}. Thin colored lines display the objective of the update step \eqref{eqn:update}. For growing $\alpha$ this illustrates how $\delta_t$ grows much slower for \NGN{} compared to \SGD{}.}
    \label{fig:toy-delta-illustration}
\end{figure}
%
%
%
\section{Experiments}\label{sec:experiments}
%
The main objective of the experiment section is to verify whether the theoretical results are reflected in the real-world performance of \SGD{}, \SPS{}, \NGN{}, and \SPP{} for different learning tasks. For simplicity, we always consider the bound $\Omega_T$ from  \cref{thm:last-iterate}. We do this comparison in a \textbf{qualitative rather than quantitative} fashion, by looking at final training loss and bound as a function of the step size~$\alpha$.
Note that for this purpose \emph{training stability} is our main interest rather than generalization properties. Hence, we usually display the training loss.
We would like to stress that it is \emph{not} the goal of our experiments to convince the reader of the superior performance of \SPS{} or \NGN{}; their performance has been studied extensively for many deep learning tasks against competitive baselines \citep{Schaipp2024a,Islamov2025,Schaipp2023}.

\paragraph{Finite-sum problems.} In all experiments, we consider the objective given as a finite sum over data points, that is $f(x) = \frac{1}{n} \sum_{i=1}^n f_i(x)$ for some $f_i:\R^d \to \R$. For batch size $b\in \N$, we can define $\mathcal{S}$ to be the set of all subsets of cardinality $b\in \N$ of $[n]$ ($s \in \mathcal{S}$ containing the indices of the batch).
Now, for $s=\{j_1,\ldots,j_b\}$ define $f(x,s) := \frac{1}{b}\sum_{i=1}^b f_{j_i}(x)$. Then, if $s$ follows a uniform distribution, we obtain $\frac{1}{n}\sum_{i=1}^n f_i(x) = f(x)= \E_s[f(x,s)]$.

\paragraph{General setup.} In all experiments, we parametrize the step size/learning rate as $\alpha_t = \alpha \cdot \eta_t$, where we call $\alpha > 0$ the \emph{base learning-rate}, and $(\eta_t)_t \subset \R_{> 0}$ is called the \emph{schedule}. Our main interest is the performance of methods as a function of $\alpha$, in particular their stability when $\alpha$ becomes large. By default, we set the schedule $(\eta_t)_t$ to be constant.
For \SPS{}, if not explicitly mentioned otherwise, we set the lower bound $C_t=0$ in every iteration.

%
\subsection{non-convex Deep Learning Tasks}
We train a \texttt{ResNet20} on \texttt{CIFAR10}. 
We run \SGD{}, \SPS{}, and \NGN{} for $20$ epochs on a range of base learning-rates $\alpha\in [10^{-2}, 10^2]$. This rather short training horizon will be sufficient to analyze stability wrt.\ large learning rates: in our experiments, the initial training stage is usually decisive for stability.
%
For \SGD{}, we run two settings:
\begin{enumerate}[label=(\arabic*)]
    \item \textbf{(without warmup):} constant learning rate schedule, that is, $\eta_t=1$ for all $t\in \N$,
    \item \textbf{(with warmup):} constant schedule after $100$ iterations of linear warmup, from $\eta_1 = 10^{-10}$ to $\eta_{100} = 1$.
\end{enumerate}

We log the gradient norm $\|g_t\|$ and batch loss $f(x_t,s_t)$ in every step, which allows us to track $\{\delta_t^{\SGD{}}, \delta_t^{\SPS{}}, \delta_t^{\NGN{}}\}$. We approximate $\E[\delta_t]$ by averaging over three different seeds (with identical initialization but different data ordering).
%

\textit{Discussion:} \cref{fig:cifar10-loss-and-bound} shows that, despite the problem at hand being non-convex, the bound from \cref{thm:last-iterate} closely reflects the range of base learning-rate for which we obtain a good final training loss. In particular, it reflects the fact that \SGD{} starts to degrade for $\alpha \geq 1$, whereas \SPS{} and \NGN{} work well up to $\alpha \approx 10$. For example, the base learning-rate for which the bound of \SPS{} is minimal, matches the one for which the \emph{actual train loss} is smallest. When using \SGD{} with warmup (\cref{fig:cifar10-loss-and-bound-warmup}), it becomes more stable and works reasonably for roughly one order of magnitude larger $\alpha$. Again, this is closely reflected in the bound as well.

For \cref{fig:cifar10-loss-and-bound} we set the unknown parameter $D = 50$ based on $\|x_T - x_1\| \approx 50$.\footnote{However, this is highly approximative, as two different methods/ learning rates can result in similar final training loss, but rather distinct final model norm.
}
and $T=20$ epochs. We ablate these choices in \cref{fig:cifar10-multipleD,fig:cifar10-loss-and-bound-ablate-T}, showing that it does not impact our main observation. Further, \cref{fig:cifar10-loss-and-bound-ablate-val} shows that the stability in terms of validation loss is very similar.

\cref{fig:cifar10-evolution-delta} depicts the evolution of $\delta_t$ over time: for \SPS{} and \NGN{}, the values of $\delta_t$ early in training grow much slower compared to \SGD{} with the same $\alpha$; this confirms our theoretical finding that $\{\delta_t^{\SPS{}}, \delta_t^{\NGN{}} \} \leq \delta_t^{\SGD{}}$ even though the learning task at hand is non-convex. 

%
% Figure with train loss and bound, w and wo warmup 
\begin{figure}[t]
    \begin{subfigure}{0.49\columnwidth}
       \centering
        \includegraphics[width=0.99\linewidth]{plots/cifar10_resnet20/stability_train_loss_epoch_19_bound.pdf}
        \caption{without warmup}
        %\label{fig:enter-label} 
    \end{subfigure}
    %
    \begin{subfigure}{0.49\columnwidth}
       \centering
        \includegraphics[width=0.99\linewidth]{plots/cifar10_resnet20/warmup/stability_train_loss_epoch_19_bound.pdf}
        \caption{with warmup}
        \label{fig:cifar10-loss-and-bound-warmup}
    \end{subfigure}
    \caption{\texttt{ResNet20} on \texttt{CIFAR10}: Actual training loss \textbf{(solid lines)} and value of the bound $\Omega_T^{\text{last}}$,~\cref{thm:last-iterate} \textbf{(dashed)} with $D=50$. After $T = 20$ epochs, \SPS{} and \NGN{} are more stable than \SGD{} for large $\alpha$ and achieve a smaller loss. This behavior is (qualitatively) reflected in the bound. \textbf{(Right)} Warmup allows \SGD{} to use a larger learning rate. Again, this is reflected in the bound $\Omega_T^{\text{last}}$.}
    \label{fig:cifar10-loss-and-bound}
\end{figure}
%
% Figureof delta evolution, w and wo warmup 
\begin{figure}[t]
    \begin{subfigure}{0.49\columnwidth}
       \centering
        \includegraphics[width=0.99\linewidth]{plots/cifar10_resnet20/deltas.pdf}
        \caption{without warmup}
        %\label{fig:enter-label} 
    \end{subfigure}
    %
    \begin{subfigure}{0.49\columnwidth}
       \centering
        \includegraphics[width=0.99\linewidth]{plots/cifar10_resnet20/warmup/deltas.pdf}
        \caption{with warmup}
        %\label{fig:enter-label} 
    \end{subfigure}
    \caption{\textbf{(Left)} $\delta_t$ explodes for \SGD{} (without warmup) in the first iterations when $\alpha$ is large. \textbf{(Right)} When using warmup over $100$ steps, the increase of $\delta_t$ is slowed down.}
    \label{fig:cifar10-evolution-delta}
\end{figure}
%


%
%
\newcommand{\bA}{\mathbf{A}}
\newcommand{\bb}{\mathbf{b}}

\subsection{Convex Regression Tasks}

We also run a set of convex tasks, where the assumptions of the theory are satisfied.
First, for ridge regression on synthetic data, we compare \SGD{}, \SPP{} and \SPS{}; we omit \NGN{} here for the sake of clarity of presentation.
We further compare \SGD{}, \SPS{} and \NGN{} on linear classification tasks.

\paragraph{Linear regression.}
%
Assume we have given a feature matrix $\bA \in \R^{n \times d}$ and targets $\bb \in \R^n$.
The objective is given by $f(x) = \frac{1}{2n} \|\bA x-\bb\|_2^2$.
For a batch $s=\{i_1,\ldots, i_b\}$, denote by $\bA_s\in \R^{b \times d}$ and $\bb_s \in \R^b$ the subset of rows of $\bA$ and $\bb$ indexed by $s$; hence, we have $f(x,s) = \frac{1}{2b} \|\bA_s x - \bb_s\|_2^2$. \cref{lem:ridge-spp-update} in the Appendix derives the \SPP{} update for this specific problem.



\textit{Discussion:} We set $n=50$, $d=10$, $b=5$. The data $\bA,~\bb$ is synthetically generated, in a way such that every datapoint can be perfectly fitted (hence we have interpolation, $\sigma^2 = 0$). The schedule is chosen to be constant for all methods.

\cref{fig:linreg-bound-and-delta} shows that \SPS{} and \SPP{} achieve reasonably good performance for values of $\alpha$ up to $10^2$, while \SGD{} starts to fail at $\alpha \approx 0.03$. Again, the bound from \cref{thm:last-iterate} closely reflects this behavior; it also matches the true performance regarding the fact that \SPP{} ``overtakes" \SPS{} for values of $\alpha \geq 1$. Next, we run three problem variations:
%
\begin{enumerate}[label=(V\arabic*)]
    \item \label{item:ablation-V1} We add noise to the targets $\bb$. As $n > d$, this results in $f_\star >0$, while $\inf_z f(z,s)=0$ still holds (as the batch size is smaller than $d$). Consequently, our theory predicts \SPS{} and \SPP{} to perform worse now.
    \item \label{item:ablation-V2} We run the noisy setting with larger batch size of $b=25$ (and 50 epochs instead of 10 to keep the number of steps equal). As $\inf_z f(z,s) >0$, we don't expect improvement for \SPS{} (as $C_t=0$), but for \SPP{} to regain advantage over \SGD{} as $\sigma^2$ is smaller than in \ref{item:ablation-V1}.
    \item \label{item:ablation-V3} We again run the noiseless case, but now with $C_t = -2$ for \SPS{}. Here, even though $\sigma^2=0$ we expect \SPS{} to loose all advantage compared to \SGD{}.
\end{enumerate}

\cref{fig:linreg-ablation} shows that for all variants the actual performance is in line with what we would expect according to our theory. 

% Figure with SPP for ridge regression
\begin{figure}[t]
    \begin{subfigure}{0.49\columnwidth}
       \centering
        \includegraphics[width=0.99\linewidth]{plots/linreg_v1/stability_train_loss_epoch_9_bound.pdf}
        \caption{Train loss and bound}
        %\label{fig:enter-label} 
    \end{subfigure}
    %
    \begin{subfigure}{0.49\columnwidth}
       \centering
        \includegraphics[width=0.99\linewidth]{plots/linreg_v1/deltas.pdf}
        \caption{Evolution of $\delta_t$}
    \end{subfigure}
    \caption{Linear regression: \textbf{(Left)} Actual training loss \textbf{(solid lines)} and value of the bound $\Omega_T^{\text{last}}$,~\cref{thm:last-iterate} \textbf{(dashed)} with $D=1$. \textbf{(Right)} The values of $\delta_t$ quickly explode for \SGD{}, whereas for \SPP{} and \SPS{} they remain bounded when $\alpha$ becomes large.}
    \label{fig:linreg-bound-and-delta}
\end{figure}
%

% Ablation for ridge regression
\begin{figure}[t]
    \begin{subfigure}{0.32\columnwidth}
       \centering
        \includegraphics[width=0.99\linewidth]{plots/linreg_v2/stability_train_loss_epoch_9_bound.pdf}
        \caption{Variant \ref{item:ablation-V1}}
    \end{subfigure}
    %
    \begin{subfigure}{0.32\columnwidth}
       \centering
        \includegraphics[width=0.99\linewidth]{plots/linreg_v4/stability_train_loss_epoch_49_bound.pdf}
        \caption{Variant \ref{item:ablation-V2}}
    \end{subfigure}
    %
    \begin{subfigure}{0.32\columnwidth}
       \centering
        \includegraphics[width=0.99\linewidth]{plots/linreg_v3/stability_train_loss_epoch_9_bound.pdf}
        \caption{Variant \ref{item:ablation-V3}}
    \end{subfigure}
    %
    \caption{Variations for the linear regression problem.}
    \label{fig:linreg-ablation}
\end{figure}
%


\paragraph{Multi-class logistic regression.}
%Given a feature matrix $\bA \in \R^{n \times d}$ and class labels $\bb \in [C]^n$, we train a linear classifier using the cross-entropy loss. 
We run multi-class logistic regression for two datasets (\texttt{vowel} and \texttt{dna}) from \texttt{LIBSVM} \citep{Chang2011}. Further dataset details are reported in \cref{sec:app:experiments-supplementary}.

\textit{Discussion:} For both datasets (see \cref{fig:vowel-bound-and-delta,fig:dna-bound-and-delta}), we again observe that the range of step sizes with good final loss roughly coincides with the range obtained from the theoretical bound (here we set $D=100$ for \texttt{vowel} and $D=50$ for \texttt{dna}). In particular, for \texttt{vowel} there is almost no advantage of \SPS{} over \SGD{}, which is likely due to the lower-bound estimation error. This confirms that even for badly specified lower bounds, \SPS{} does not perform worse than \SGD{}, but it also shows no advantage. \NGN{} is affected by this to a lesser extent.

% Figure for vowel
\begin{figure}[H]
    \begin{subfigure}{0.49\columnwidth}
       \centering
        \includegraphics[width=0.99\linewidth]{plots/vowel/stability_train_loss_epoch_19_bound.pdf}
        \caption{Train loss and bound}
        %\label{fig:enter-label} 
    \end{subfigure}
    %
    \begin{subfigure}{0.49\columnwidth}
       \centering
        \includegraphics[width=0.99\linewidth]{plots/vowel/deltas.pdf}
        \caption{Evolution of $\delta_t$}
    \end{subfigure}
    \caption{Logistic regression on \texttt{vowel} dataset.}
    \label{fig:vowel-bound-and-delta}
\end{figure}
%

\section{Conclusion}

We have presented a theoretical analysis for step-size stability of several stochastic optimization methods. Within our framework, we prove that adaptive methods such as \SPS{} or \NGN{} are always more stable than \SGD{} in terms of a (weakly) convex convergence bound. Our experiments show a close relationship between the convex bound and the actual training loss, even for non-convex deep learning tasks.

Finally, we want to point out again the main limitations and possible directions for future work.
Our current analysis can be converted into convergence rates by applying a standard Lipschitz bound on $\delta_t$ (see \cref{sec:app:conv-comparison} for details); however, this is too coarse to capture the advantage of adaptive methods like \SPS{} and \NGN{} and hence not sufficient to show a speedup. Tightening this conversion from bound to convergence rate would be an interesting extension.

On the practical side, an open question is how tracking the stability index $\delta_t$ could be leveraged in practice, for example, to terminate unstable runs early. This paper does not explore such heuristics. Another limitation of our current work is that the analysis is restricted to \SGD-based updates, not including momentum, weight decay or preconditioning. For the Polyak step size, it has been shown that it can be combined with all of these \citep{Schaipp2024a}, and additionally recent works show that it is beneficial in combination with base optimizers like \texttt{Muon} \citep{Crawshaw2025} and \texttt{ScheduleFree} \citep{Defazio2026}. Hence, a direction for future work would be to extend our analysis to other update geometries beyond the model-based framework presented here.

%%%%%%%%%%%%%%%%%%%%%%%%%%%%%%%%%%%%%%%%%%%%%%%%%%%%%%%%%%%%%%%%%%%%%%%%%%%%%%
%%%%%%%%%%%%%%%%%%%%%%%%%%%%%%%%%%%%%%%%%%%%%%%%%%%%%%%%%%%%%%%%%%%%%%%%%%%%%%
%%%%%%%%%%%%%%%%%%%%%%%%%%%%%%%%%%%%%%%%%%%%%%%%%%%%%%%%%%%%%%%%%%%%%%%%%%%%%%
%%%%%%%%%%%%%%%%%%%%%%%%%%%%%%%%%%%%%%%%%%%%%%%%%%%%%%%%%%%%%%%%%%%%%%%%%%%%%%
%%%%%%%%%%%%%%%%%%%%%%%%%%%%%%%%%%%%%%%%%%%%%%%%%%%%%%%%%%%%%%%%%%%%%%%%%%%%%%
%%%%%%%%%%%%%%%%%%%%%%%%%%%%%%%%%%%%%%%%%%%%%%%%%%%%%%%%%%%%%%%%%%%%%%%%%%%%%%
% HERE THE 8 PAGES END
\clearpage


% Acknowledgements should only appear in the accepted version.
\section*{Acknowledgements}
We thank Umut \c{S}im\c{s}ekli and Francis Bach for comments and feedback.
A.~Taylor is supported by the European Union (ERC grant CASPER 101162889). Fabian Schaipp is supported by the French government under the management of Agence Nationale de la Recherche as part of the “Investissements d’avenir” program,
reference ANR-19-P3IA-0001 (PRAIRIE 3IA Institute), and the
European Research Council Starting Grant DYNASTY – 101039676.

\section*{Impact Statement}

This paper presents work whose goal is to advance the field of Machine
Learning. There are many potential societal consequences of our work, none
which we feel must be specifically highlighted here.

%%%%%%%%%%%%%%%%%%%%%%%%%%%%%%%%%%%%%%%%%%%%%%%%%%%%%%%%%%%%%%%%%%%%%%%%%%%%%%
%%%%%%%% References
\bibliography{lib}
\bibliographystyle{icml2026}
%%%%%%%%%%%%%%%%%%%%%%%%%%%%%%%%%%%%%%%%%%%%%%%%%%%%%%%%%%%%%%%%%%%%%%%%%%%%%%

%%%%%%%%%%%%%%%%%%%%%%%%%%%%%%%%%%%%%%%%%%%%%%%%%%%%%%%%%%%%%%%%%%%%%%%%%%%%%%%
%%%%%%%%%%%%%%%%%%%%%%%%%%%%%%%%%%%%%%%%%%%%%%%%%%%%%%%%%%%%%%%%%%%%%%%%%%%%%%%
% APPENDIX
%%%%%%%%%%%%%%%%%%%%%%%%%%%%%%%%%%%%%%%%%%%%%%%%%%%%%%%%%%%%%%%%%%%%%%%%%%%%%%%
%%%%%%%%%%%%%%%%%%%%%%%%%%%%%%%%%%%%%%%%%%%%%%%%%%%%%%%%%%%%%%%%%%%%%%%%%%%%%%%
\newpage
\appendix
\onecolumn

\section{Auxiliary Lemmas and Missing Proofs}\label{sec:app:proofs}

%%%%%%%%%%%%%%%%%%%%%%%%%%%%%%%%%%%%%%%%%%%%%%%%%%%%%%%%%%%%%%%%%%%%%%%%%%
%%% OLD result on squaring a convex fucntion
%%%%%%%%%%%%%%%%%%%%%%%%%%%%%%%%%%%%%%%%%%%%%%%%%%%%%%%%%%%%%%%%%%%%%%%%%%
% \begin{lemma}[Squaring a nonnegative convex function preserves convexity]
% \label{lem:squareconv}
% Let $C \subseteq \mathbb{R}^d$ be a convex set and let $h:C\to[0,\infty)$ be a convex function.
% Define $f:C\to\mathbb{R}$ by
% \[
% f(x) := h(x)^2.
% \]
% Then $f$ is convex on $C$.
% \end{lemma}

% \begin{proof}
% Fix any $x,y\in C$ and any $\lambda\in[0,1]$. Since $C$ is convex, the point
% \[
% z := \lambda x + (1-\lambda)y
% \]
% belongs to $C$. By convexity of $h$,
% \begin{equation}\label{eq:convex_h}
% h(z) \le \lambda h(x) + (1-\lambda)h(y).
% \end{equation}
% Because $h$ is nonnegative, we have $h(z)\ge 0$, $h(x)\ge 0$, and $h(y)\ge 0$; hence the right-hand side of \eqref{eq:convex_h} is also nonnegative. The map $\phi:[0,\infty)\to[0,\infty)$ defined by $\phi(t)=t^2$ is nondecreasing on $[0,\infty)$, so applying $\phi$ to both sides of \eqref{eq:convex_h} preserves the inequality:
% \begin{equation}\label{eq:square_monotone}
% h(z)^2 \le \big(\lambda h(x) + (1-\lambda)h(y)\big)^2.
% \end{equation}
% Next, since $\phi(t)=t^2$ is convex on $\mathbb{R}$ (equivalently on $[0,\infty)$), Jensen's inequality yields
% \begin{equation}\label{eq:jensen_square}
% \big(\lambda h(x) + (1-\lambda)h(y)\big)^2 \le \lambda h(x)^2 + (1-\lambda)h(y)^2.
% \end{equation}
% Combining \eqref{eq:square_monotone} and \eqref{eq:jensen_square} gives
% \[
% h(\lambda x+(1-\lambda)y)^2 \le \lambda h(x)^2 + (1-\lambda)h(y)^2,
% \]
% which is exactly
% \[
% f(\lambda x+(1-\lambda)y) \le \lambda f(x) + (1-\lambda)f(y).
% \]
% Since $x,y\in C$ and $\lambda\in[0,1]$ were arbitrary, $f$ is convex on $C$.
% \end{proof}
%%% END: OLD
%%%%%%%%%%%%%%%%%%%%%%%%%%%%%%%%%%%%%%%%%%%%%%%%%%%%%%%%%%%%%%%%%%%%%%%%%%


%%%%%%%%%%%%%%%%%%%%%%%%%%%%%%%%%%%%%%%%%%%%%%%%%%%%%%%%%%%%%%%%%%%%%%%%%%







\begin{lemma}\label[lemma]{lem:exp-of-min}
    Let $X, Y$ be random variables. Then $\E[\min\{X,Y\}] \leq \min\{\E[X], \E[Y]\}$.
\end{lemma}
\begin{proof}
    Follows from Jensen's inequality and the fact that $(x,y)\mapsto \min\{x,y\}$ is concave.
\end{proof}

%
\begin{lemma}\label[lemma]{lem:ridge-spp-update}
    For $\lambda > 0$, the solution to
    \begin{align*}
        x_{t+1} = \argmin_{y \in \R^d} \tfrac{1}{2b} \|\bA_{s_t} y - \bb_{s_t}\|^2 + \tfrac{\lambda}{2}\|y\|^2 + \frac{1}{2\alpha_t}\|y-x_t\|^2
    \end{align*}
    is given by the solution to the linear system
    \begin{align*}
        \big[\tfrac1b \bA_{s_t}^T \bA_{s_t} + (\lambda + \alpha_t^{-1})\Id{} \big]x_{t+1} = \alpha_t^{-1}x_t + \tfrac{1}{b} \bA_{s_t}^T \bb_{s_t}.
    \end{align*}
\end{lemma}
%
\begin{proof}
    Follows from the necessary and sufficient optimality conditions.
\end{proof}

%%%%%%%%%%%%%%%%%%%%%%%%%%%%%%%%%%%%%%%%%%%%%%%%%%%%%%%%%%%%%%%%%%%%%%%%%%
\subsection{Proofs for \cref{sec:theory-weakly-convex}}\label{sec:app:proofs-weakly-convex}

\begin{proof}[Proof of \cref{lem:weakly-convex-base}]
    Let us define $\Psi_t(x) = f_{x_t}(x,s_t) + \frac{1}{2\alpha_t}\|x-x_t\|^2$. Then, due to \ref{item:tA1}, $\Psi_t$ is $(\frac{1}{\alpha_t} - \lambda)$-strongly convex. Further, $x_{t+1}$ is a minimizer of $\Psi_t$. Therefore, for any $u\in \R^d$ it holds
    \begin{align*}
        \Psi_t(x_{t+1}) + \frac{1-\lambda\alpha_t}{2\alpha_t}\|u-x_{t+1}\|^2 \leq \Psi_t(u).
    \end{align*}
    Plugging in $u=\hat{x}_t$ and the definition of $\Psi_t$, we get
    \begin{align}\label{eqn:weak-conv-start-ineq}
        \frac{1-\lambda\alpha_t}{2\alpha_t}\|x_{t+1} - \hat{x}_t\|^2 \leq
        \frac{1}{2\alpha_t}\|\hat{x}_t - x_t\|^2 - \frac{1}{2\alpha_t}\|x_{t+1} - x_t\|^2 - [f_{x_t}(x_{t+1},s_t) - f_{x_t}(\hat{x}_t,s_t)].
    \end{align}
    Similar to the proof of \cref{lem:base}, from \ref{item:tA1} and \ref{item:tA2} we get
    \begin{align*}
        f_{x_t}(x_{t+1},s_t) - f_{x_t}(\hat{x}_t,s_t) &= 
        \underbracket{f_{x_t}(x_{t+1},s_t) - f(x_t,s_t)}_{= -\delta_t - \frac{1}{2\alpha_t}\|x_{t+1}-x_t\|^2} + f(x_{t},s_t) - f_{x_t}(\hat{x}_t,s_t) \\
        &\geq
        f(x_t,s_t) - f(\hat{x}_t, s_t) - \frac{\rho_{s_t}}{2}\|x_t-\hat{x}_t\|^2
        -\delta_t - \frac{1}{2\alpha_t}\|x_{t+1}-x_t\|^2.
    \end{align*}
    Plugging this back in \eqref{eqn:weak-conv-start-ineq},  multiplying by $2\alpha_t$, and taking conditional expectation, we get
    \begin{align*}
        (1-\lambda\alpha_t)\E_t\|x_{t+1} - \hat{x}_t\|^2 \leq
        (1+\alpha_t\rho)\|\hat{x}_t - x_t\|^2  - 2\alpha_t[f(x_{t}) - f(\hat{x}_t)] + 2\alpha_t \E_t[\delta_t].
    \end{align*}
    Since \ref{item:tA1}-\ref{item:tA2} imply weak convexity of $f$ and $\bar{\rho} > \rho+\lambda$, by definition of the proximal operator we have
    \begin{align*}
         f(x_t) - f(\hat{x}_t) \geq \frac{\bar{\rho}}{2}\|x_t-\hat{x}_t\|^2.
    \end{align*}
    Therefore, we get
    \begin{align*}
        (1-\lambda\alpha_t)\E_t\|x_{t+1} - \hat{x}_t\|^2 \leq
        \big(1 - \alpha_t(\bar{\rho}-\rho)\big)\|x_t - \hat{x}_t\|^2
        + 2\alpha_t \E_t[\delta_t].
    \end{align*}
    Multiplying by $(1-\lambda\alpha_t)^{-1}$ gives the final result.
\end{proof}
%


%
\begin{proof}[Proof of \cref{thm:weakly-convex}]
    Denote again $\hat{x}_t=\prox{(1/\bar{\rho})f}(x_t)$. By definition of the Moreau envelope and applying \cref{lem:weakly-convex-base}, we have
    \begin{align*}
        \E_t[\env{f}^{1/\bar{\rho}}(x_{t+1})] &\leq
        \E_t[f(\hat{x}_t) + \frac{\bar{\rho}}{2}\|\hat{x}_t - x_{t+1}\|^2] \\
        &\leq f(\hat{x}_t) +
        \frac{\bar{\rho}}{2} \Big[\|x_{t} - \hat{x}_t\|^2 - \frac{\alpha_t(\bar{\rho}-\rho-\lambda)}{1-\alpha_t\lambda}\|x_{t} - \hat{x}_t\|^2 + \frac{2\alpha_t}{1-\alpha_t\lambda} \E_t[\delta_t]\Big].
    \end{align*}
For the third term, we use that $\|x_{t} - \hat{x}_t\|^2 = (\bar{\rho})^{-2}\|\nabla \env{f}^{1/\bar{\rho}}(x_t)\|^2$ \citep[Lem.\ 2.2]{Davis2019}. Thus, we get
\begin{align*}
    \E_t[\env{f}^{1/\bar{\rho}}(x_{t+1})] \leq \env{f}^{1/\bar{\rho}}(x_{t})
    - \frac{\alpha_t(\bar{\rho}-\rho-\lambda)}{2\bar{\rho}(1-\alpha_t\lambda)} \|\nabla \env{f}^{1/\bar{\rho}}(x_t)\|^2
    + \frac{\alpha_t\bar{\rho}}{1-\alpha_t\lambda} \E_t[\delta_t].
\end{align*}
Applying total expectation, and summing over $t=k,\ldots,T$, we obtain
\begin{align*}
    \sum_{t=k}^{T} \frac{\alpha_t}{1-\alpha_t\lambda} \E\|\nabla \env{f}^{1/\bar{\rho}}(x_t)\|^2
    \leq \frac{2\bar{\rho}}{\bar{\rho}-\rho-\lambda} \big[\E[\env{f}^{1/\bar{\rho}}(x_{k})] - \E[\env{f}^{1/\bar{\rho}}(x_{T+1})]\big] + \frac{2\bar{\rho}^2}{\bar{\rho}-\rho-\lambda}\sum_{t=k}^{T}\frac{\alpha_t}{1-\alpha_t \lambda} \E[\delta_t].
\end{align*}
\end{proof}

%%%%%%%%%%%%%%%%%%%%%%%%%%%%%%%%%%%%%%%%%%%%%%%%%%%%%%%%%%%%%%%%%%%%%%%%%%

\subsection{Proofs for \cref{sec:models}}\label{sec:app:proofs-deltas}

For the following proofs, recall the notation from \eqref{eqn:psi-delta}, that is
\begin{align*}
\Psi_t=f_{x_t}(x_{t+1},s_t) + \frac{1}{2\alpha_t}\|x_{t+1}-x_t\|^2.
\end{align*}

%%%%%RESTATE LEMMA%%%%%%%%%
%\lemdeltasps*
%%%%%%%%%%%%%%%%%%%%%%%%%%
\begin{proof}[Proof of \cref{lem:delta-sps}]
Denote $C_t:=C_{s_t}$.
Due to \citet[Lem.\ 9]{Schaipp2023} it holds $x_{t+1} = x_t - \tau_t g_t$ with
\begin{align*}
    \begin{split}
    %\tau_t &:= \min\{\alpha_t, \frac{f(x_t,s_t)-C_{t}}{\|g_t\|^2}\} \geq 0,\\
    \Psi_t &= f(x_t, s_t) - \tau_t\|g_t\|^2 + \frac{\tau_t^2}{2\alpha_t}\|g_t\|^2 \\
    &= f(x_t, s_t) - \tau_t[1-\frac{\tau_t}{2\alpha_t}]\|g_t\|^2.
    \end{split}
\end{align*}
Since $ \Psi_t = f(x_t,s_t) -\delta_t^{\SPS{}}$  we have from the above that $\delta_t^{\SPS{}} = \tau_t[1-\frac{\tau_t}{2\alpha_t}]\|g_t\|^2$.
Given that the function $\tau \mapsto \tau(1-\frac{\tau}{2\alpha})$ is maximized at $\hat \tau=\alpha$, we conclude that $\delta_t^{\SPS{}} \leq \frac{\alpha_t}{2}\|g_t\|^2 = \delta_t^{\SGD{}}$ (assuming both $\delta_t$ would be measured at the same iterate $x_t$ and batch $s_t$).

Now, we treat the two cases separately:
\begin{itemize}
    \item If $\tau_t = \frac{f(x_t,s_t)-C_t}{\|g_t\|^2}$, we have $1-\frac{\tau_t}{2\alpha_t} \leq 1$ and hence $\delta_t^{\SPS{}} \leq \tau_t \|g_t\|^2 = f(x_t,s_t)-C_t \leq \alpha_t \|g_t\|^2$.
    \item In the other case, if $\tau_t=\alpha_t$ then $\delta_t^{\SPS{}} = \tau_t[1-\frac{\tau_t}{2\alpha_t}]\|g_t\|^2 = \frac{\alpha_t}{2}\|g_t\|^2 \leq \frac{1}{2}[f(x_t,s_t)-C_t]$.
\end{itemize}
%
Combining both cases we can upper-bound 
\begin{equation}
    \delta_t^{\SPS{}} \leq \min\{\alpha_t\|g_t\|^2, f(x_t,s_t)-C_t \}.
\end{equation}
Consequently $\delta_t^{\SPS{}}$ scales sub-linearly with $\alpha_t$.
\end{proof}


%%%%%%%%RESTATE LEMMA%%%%%%%%
%\lemdeltangn*
%%%%%%%%%%%%%%%%%
\begin{proof}[Proof of \cref{lem:delta-ngn}]
For $x_t\in \R^d$ and $g_t\in \partial f(x_t,s_t)$, denote $f_t:=f(x_t,s_t)$. \citet{Orvieto2024} show that update \eqref{eqn:update} with model \eqref{eqn:ngn-model} is given by \eqref{eq:NGNstep}.
% Let us compute:
% \begin{align*}
%     f_{x_t}(x_{t+1},s_t) &= \Big(\sqrt{f_t} - \frac{1}{2\sqrt{f_t}}\gamma_t \|g_t\|^2\Big)^2 = f_t - \gamma_t \|g_t\|^2 + \frac{1}{4f_t}\gamma_t^2 \|g_t\|^4
%     ,
% \end{align*}
Plugging into \eqref{eqn:ngn-model}, we get
\begin{align*}
    f_{x_t}(x_{t+1},s_t) &= 
    %\Big(\sqrt{f_t} - \frac{1}{2\sqrt{f_t}}\gamma_t \|g_t\|^2\Big)^2 =
    f_t - \gamma_t \|g_t\|^2 + \frac{1}{4f_t}\gamma_t^2 \|g_t\|^4,
\end{align*}
and $\frac{1}{2\alpha_t}\|x_{t+1}-x_t\|^2 = \frac{1}{2\alpha_t}\gamma_t^2\|g_t\|^2$.
Hence,
\begin{align*}
    %\hspace{-5ex}f_{x_t}(x_{t+1},s_t) + \frac{1}{2\alpha_t}\|x_{t+1}-x_t\|^2
    \hspace{-3ex}\Psi_t
    &= f_t - \gamma_t \|g_t\|^2 + \gamma_t^2 \underbracket{\big[\frac{1}{4f_t} \|g_t\|^4 + \frac{1}{2\alpha_t}\|g_t\|^2 \big]}_{= \frac{\|g_t\|^2}{2\alpha_t}(1+\frac{\alpha_t\|g_t\|^2}{2f_t}) = \frac{\|g_t\|^2}{2\gamma_t}}\\
    &= f_t - \frac{\gamma_t}{2} \|g_t\|^2.
\end{align*}
Altogether, we obtain $\delta_t^{\NGN{}} = f_t - \Psi_t = \frac{\gamma_t}{2} \|g_t\|^2.$ From $\gamma_t \leq \alpha_t$ it follows that $\delta_t^{\NGN{}} \leq \delta_t^{\SGD{}}$. 
\end{proof}
%%%%%%%%%%%%%%%%%%%%%%%%%%%%


%%%%%%%%%%%%%%%%
%\lemdeltaspp*
%%%%%%%%%%%%%%%%
\begin{proof}[Proof of \cref{lem:delta-spp}]
    We have
    $\Psi_t  - f(x_t,s_t) = f(x_{t+1}, s_t) - f(x_t,s_t) + \frac{1}{2\alpha_t}\|x_{t+1}-x_t\|^2.$
    For any $\alpha_t \in (0, 1/\bar{\eta})$, it holds
    \begin{align*}
        \Psi_t = \min_{y\in \R^d} f(y,s_t) + \frac{1}{2\alpha_t}\|y-x_t\|^2
        \geq \min_{y\in \R^d} f(y,s_t).
    \end{align*}
    Hence, we have
    \begin{align}\label{eq:spp-bound-one}
    \begin{split}
        \Psi_t  - f(x_t,s_t) &\geq \inf_{y\in \R^d} f(y,s_t) - f(x_t,s_t).
        %\\&= - \big[f(x_t,s_t) - f(x_\star, s_t) + f(x_\star, s_t) - \inf_z f(z,s_t) \big].
    \end{split}
    \end{align}
    On the other hand, due to $\bar{\eta}$-weak convexity of $f(\cdot,s)$ we have
    \begin{align*}
        & \hspace{-3ex} f(x_{t+1},s_t) \geq f(x_t,s_t) + \iprod{g_t}{x_{t+1}-x_t} - \frac{\bar{\eta}}{2}\|x_{t+1}-x_t\|^2 \\
        &\geq f(x_t,s_t) - \frac{\epsilon}{2}\|g_t\|^2 - (\frac{1}{2\epsilon}+ \frac{\bar{\eta}}{2} )\|x_{t+1}-x_t\|^2,
    \end{align*}
    where we use Young's inequality with $\epsilon>0$ in the second step. Setting $\epsilon=\frac{\alpha_t}{1-\alpha_t \bar{\eta}} > 0$, we get
    \begin{align}\label{eq:spp-bound-two}
        \Psi_t - f(x_t,s_t) \geq -\frac{\alpha_t}{2(1-\alpha_t \bar{\eta})}\|g_t\|^2.
    \end{align}
    Combining \eqref{eq:spp-bound-one} and \eqref{eq:spp-bound-two} , we have
    \begin{align*}
        &\hspace{-3ex}\Psi_t - f(x_t,s_t) \\
        & \geq \max\{-\frac{\alpha_t}{2(1-\alpha_t \bar{\eta})}\|g_t\|^2 ,\inf_{y\in \R^d} f(y,s_t) - f(x_t,s_t)\} \\
        &= -\min\{\frac{\alpha_t}{2(1-\alpha_t \bar{\eta})}\|g_t\|^2, f(x_t,s_t) - \inf_{y\in \R^d} f(y,s_t)\}.
    \end{align*}
    From \eqref{eqn:psi-delta}, this implies~$\delta_t^{\SPP{}} \leq \min\{\frac{\alpha_t}{2(1-\alpha_t \bar{\eta})}\|g_t\|^2, f(x_t,s_t) - \inf_{y\in \R^d} f(y,s_t)\}$.
\end{proof}

%%%

\subsection{Quadratics as a Special Case}\label{sec:app:quadratics}

Consider the class of functions $f(x,s):=\frac{s^2}{2}\|x-x_\star\|^2$ for $s\in \R$ and some fixed $x_\star \in \R^d$. This represents a linear regression problem with interpolation condition, i.e.\ $x_\star = \argmin_x f(x,s)$ for all $s\in \mathcal{S}$ and hence $\sigma^2=0$. In this case, the formulae for  $\delta_t^{\{\texttt{SGD},\texttt{SPS},\texttt{NGN},\texttt{SPP}\}}$ simplify a lot, and therefore this serves as a good illustration of the different stability behaviors.

First, note that $\|\nabla f(x,s)\|^2 = 2s_t^2f(x,s)$ for any $x\in\R^d,~s\in\R$. From now on, assume that $x_t\in \R^d$ is the (fixed) current iterate, and denote $f_t:=f(x_t,s_t) = \frac{\|g_t\|^2}{2s_t^2}$. For \texttt{SPS},  we can clearly use $C_t=0$ in this case, and obtain $\tau_t = \min\{\alpha_t, \frac{f_t}{\|g_t\|^2}\} = \min\{\alpha_t, \frac{1}{2s_t^2}\}$. Plugging into the exact formula for $\delta_t^{\texttt{SPS}}$ from \cref{lem:delta-sps}, and doing a simple case distinction, we obtain $\delta_t^{\texttt{SPS}} \leq \min\{\frac{\alpha_t}{2}, \frac{1}{2s_t^2}\}\|g_t\|^2$. Note that this bound is relatively tight, and can be made exact at the cost of a less compact expression.
%
For \NGN{}, we have $\gamma_t = \frac{\alpha_t}{1+\alpha_ts_t^2} \leq \frac{1}{s_t^2}$ for all $\alpha_t > 0$.
%
For \SPP{}, we can compute $x_{t+1} = \frac{1}{1+\alpha_t s_t^2}[x_t + \alpha_ts_t^2x_\star]$ and $\frac{1}{2\alpha_t}\|x_{t+1}-x_t\|^2 = \frac12 \frac{\alpha_ts_t^4}{(1+\alpha_ts_t^2)^2}\|x_t-x_\star\|^2 =  \frac{\alpha_t s_t^2}{(1+\alpha_ts_t^2)^2}f_t$. This leads to $\delta_t^{\texttt{SPP}} = \frac{\alpha_ts_t^2}{1+\alpha_ts_t^2}f_t$.


Altogether, we obtain the following stability indices:
\begin{align*}
    \delta_t^{\texttt{SGD}} &= \frac{\alpha_t}{2}\|g_t\|^2, \\
    \delta_t^{\texttt{SPS}} &\leq \min\{\frac{\alpha_t}{2}, \frac{1}{2s_t^2}\}\|g_t\|^2, \\
    \delta_t^{\texttt{NGN}} &= \frac{\alpha_t}{2(1+\alpha_ts_t^2)}\|g_t\|^2, \\
    \delta_t^{\texttt{SPP}} &= \frac{\alpha_t}{2}\frac{1}{1+\alpha_ts_t^2}\|g_t\|^2. \\
\end{align*}
In conclusion, only $\delta_t^{\texttt{SGD}}$ scales linearly with $\alpha_t$, whereas for all other methods, their stability index does not grow linearly when $\alpha_t\to \infty$.

\section{Discussion of Related Work}

\subsection{Comparison of Convergence Theorems to the Literature} \label{sec:app:conv-comparison}

Here we consider our resulting convergence results \cref{thm:avg-iterate,thm:last-iterate}, when instantiated with each model, and compare to what is known in the literature. For the linear model and the resulting \SGD{} method, we recover the standard rates as pointed out in~\cref{sec:linearmodel}.

For \eqref{eqn:sps} and \eqref{eqn:ngn} we have shown in~\cref{lem:delta-sps} and \cref{lem:delta-ngn} respectively, that these methods have a stability constant $\delta_t$ that is less or equal than that of of \SGD{}, that is 
\begin{equation} \label{eq:zefzefzfzeg}
    \delta_t \;\leq \; \delta_t^{\SGD{}} =  \frac{\alpha_t}{2}\|g_t\|^2.
\end{equation}  Using this bound in \cref{thm:avg-iterate,thm:last-iterate} we immediately arrive at a convergence result for \SPS{} and \NGN{}. However, this convergence result will be pessimistic, since we bound away the possible stability benefits of \SPS{} and \NGN{} when using~\eqref{eq:zefzefzfzeg}. 

Using the pessimistic upper bound~\eqref{eq:zefzefzfzeg}
in~\cref{thm:avg-iterate} with a constant step size $\alpha_t \equiv \alpha$ we arrive at the convergence rate
    \begin{align} \label{eq:theospspess}
        \bar f_T - f(x_\star) \leq \frac{D^2}{2T \alpha } + \alpha \frac{\sum_{t=1}^T  \E[\|g_t\|^2]}{2T}.
    \end{align}
    By further assuming the loss is Lipschitz continuous with $\|g_t\| \leq G$ and a step size of $\alpha = \frac{1}{\sqrt{T}}$ this gives a $\mathcal{O}(1/\sqrt{T})$ complexity.
This is exactly the rate of \SGD{} in the convex and Lipschitz setting, 
see for the an early such proof~\citet{pmlr-v28-shamir13} and Theorem 9.7 in~\citet{handbook} for a didactic reference.

For~\ref{eqn:sps} (introduced by~\citet{Loizou2021} and called \SPS{$_{\max}$} therein), as far as we are aware, the only result for convex and Lipschitz losses also requires interpolation \citep[Thm.\ C.1]{Loizou2021}, which gives a $\mathcal{O}(1/\sqrt{T})$ complexity.

As for the \NGN{} method, there currently exists only an analysis of the method for the Lipschitz-smooth setting. Thus~\eqref{eq:theospspess} is the first complexity for \NGN{} in the convex non-smooth setting.

%Finally for the \SPP{} method, our result~\eqref{eq:theospspess} matches that of~\citet[Thm.\ 4.1]{Davis2019} under the same assumptions.

\subsection{Additional Remarks on Related Work}\label{sec:app:related-work}
%
The theory by \citet{Loizou2021} suggests that \SPS{} can achieve the standard convergence rates without knowledge of the Lipschitz(-smoothness) constant.
More concretely, \citet[Thm.\ 3.4]{Loizou2021} shows that if $\alpha_t = \alpha$ for all $t\in \N$, and if each $f(\cdot,s)$ is convex and $L_s$-smooth, we have
\begin{align*}
    \E[f(\frac{1}{T}\sum_{t=1}^T x_t) - f(x_\star)] \leq \frac{\|x_1-x_\star\|^2}{\rho T} + \frac{2\sigma^2 \alpha}{\rho}.
\end{align*}
Here $\rho := \min\{\frac{1}{2L_{\max}}, \alpha\}$ and $L_{\max} = \max_s L_s$.
However, this theory can not explain why \SPS{} is less sensitive to the choice of $\alpha$: for $\alpha \geq \frac{1}{2L_{\max}}$, the bound is increasing in $\alpha$ linearly (as $\rho=\frac{1}{2L_{\max}}$). Moreover, in the smooth case, it holds 
$$\frac{f(x_t,s_t)-\inf_z f(z, s_t)}{\|g_t\|^2} \geq \frac{1}{2L_{s_t}} \geq \frac{1}{2L_{\max}}.$$
For \texttt{ResNet20} on \texttt{CIFAR10}, the term on the left is around $10^{-1}$ in early training, suggesting that $2 L_{\max} \geq 10^1$ and hence the bound will be increasing for $\alpha \geq 10^{-1}$. But in practice, \SPS{} achieves the best performance for much larger values of $\alpha$ (in between $1$ and $10$, see \cref{fig:cifar10-loss-and-bound}).

For stochastic proximal point (\SPP{}), several analyses show that convergence can be proven for any step size, for example \citet[Thm.\ 5.3, 6.4]{Tovmasyan2025}; however, these analyses derive rates for \SPP{} under various set of (additional) assumptions, such as smoothness, strong convexity, or interpolation; our analysis on the other hand focuses on the one-step stability index $\delta_t$, which allows to compare stability across methods.

\citet{Orabona2026} show negative results for stability and convergence of the Polyak step size; however, their results are not in conflict with ours, as their examples involve suboptimal estimates for $C_s$ and require large values for $\alpha_t$ (some results only hold for $\alpha_t = + \infty$). This is in line with our analysis, as the stability index becomes worse when both $\epsilon_{\text{lb}}$ and $\alpha_t$ grow simultaneously.

\section{Experiments: Supplementary Material}\label{sec:app:experiments-supplementary}

\paragraph{Dataset generation for linear regression.}
The data generation procedure was adapted from \citet{Loizou2021} and the reference therein: we first generate a matrix $\bA \in\R^{n\times d}$ with entries from a standard normal. We add one to all entries to introduce a bias. We then scale each column with a factor drawn from $\mathcal{N}(0,1)$ multiplied by ten. We now sparsify $\bA$ by setting each entry to zero with probability $1-30\tfrac{\ln(n)}{n}$. We finally scale each column to have norm $10$ (this step affects mainly the scaling of the step size).

We then generate a vector $\hat{x} \in\R^d$ with entries from a standard normal, re-scaled such that $\|\hat{x}\|=1$. We generate $\bb = \bA \hat{x}$; in the noisy setting, we add a vector with entries from a standard normal to $\bA$. 


\section{Generalizations of the \NGN{} Step Size}\label{sec:app:ngn-generalizations}

The \NGN{} model \eqref{eqn:ngn-model} is constructed as square-root $\to$ linearize $\to$ square. This idea can be generalized to other functions on the non-negative halfspace, and their respective inverse function. Below, we derive the model-based update when using instead $\log$ $\to$ linearize $\to$ $\exp$. We will see that the resulting update is related to the Lambert--$W$ function.

\paragraph{General maps.} Let $U\subseteq  \R_{\geq 0}$ and $\varphi: U \to \R$ and $\psi: \R \to \R$ satisfy the following:
\begin{enumerate}[label=(B\arabic*)]
    \item \label{assum:phi-psi-one} For any $x \in U$ we have $\psi(\varphi(x)) = x$.
    \item \label{assum:phi-psi-two} $\psi$ is differentiable and convex on $\R$.
    \item \label{assum:phi-psi-three} $\varphi$ is differentiable on $U$ with $\varphi'(x) \neq 0$ for all $x\in U$.
\end{enumerate}
Second, \ref{assum:phi-psi-one} implies that $\varphi$ is injective. Therefore, $\varphi$ must be either monotonically increasing or decreasing. If $\varphi$ is concave, then its first derivative is non-increasing, and hence $\varphi'(0) >0$ is sufficient for \ref{assum:phi-psi-three} to hold with $U=\R_{\geq 0}$. 

We will usually choose $U=\R_{>0}$ or $U=\R_{\geq 0}$ depending on the loss function being strictly positive or non-negative. (In the latter case, assume that there exists a continuous extension of $\varphi'$ on $\R_{\geq 0}$ with $\varphi'(0)\neq 0$).  

Then, for $g\in \partial f(x,s)$, define the model
\begin{align} \label{eq:phi-psi-model}
    f_x(y,s) = \psi\big(\varphi(f(x,s)) + \varphi'(f(x,s))\iprod{g}{y-x}\big).
\end{align}

First, note that Assumption~\ref{item:A1} holds due to convexity of $\psi$. Denote $f_t := f(x_t,s_t)$, and 
let $u_t = \varphi'(f_t)$. It holds
\begin{align*}
    x_{t+1} = \argmin_{y\in \R^d}  f_{x_t}(y,s_t) + \frac{1}{2\alpha_t} \|y-x_t\|^2
\end{align*}
if and only if 
\begin{align*}
    x_{t+1} = x_t - \alpha_t u_t v_t g_t, \quad v_t := \psi'\big(\varphi(f_t) + \varphi'(f_t)\iprod{g_t}{x_{t+1}-x_t}\big). 
\end{align*}
Therefore, we get the condition
\begin{align*}
    v_t = \psi'\big(\varphi(f_t) - \alpha_t v_t u_t^2 \|g_t\|^2\big).
\end{align*}
We can compute
\begin{align*}
    f_{x_t}(x_{t+1},s_t) = \psi\big(\varphi(f_t) - \alpha_t v_t u_t^2 \|g_t\|^2\big),
\end{align*}
and hence
\begin{align}\label{eqn:psi-phi-delta}
\begin{split}
    \delta_t &= f_t - \frac{1}{2\alpha_t} \|x_{t+1}-x_t\|^2 - f_{x_t}(x_{t+1},s_t)  \\
    &= f_t - \frac12 \alpha_t u_t^2 v_t^2 \|g_t\|^2 - \psi\big(\varphi(f_t) - \alpha_t v_t u_t^2 \|g_t\|^2\big).
\end{split}
\end{align}
Let us focus on the last term: using convexity, we have $\psi(a+b) \geq \psi(a) + \psi'(a)b$, and hence
\begin{align*}
    \psi\big(\varphi(f_t) - \alpha_t v_t u_t^2 \|g_t\|^2\big) \geq \psi(\varphi(f_t)) - \psi'(\varphi(f_t)) \alpha_t v_t u_t^2 \|g_t\|^2.
\end{align*}
%

Now, from \ref{assum:phi-psi-one} and the chain rule it follows that $1 = \psi'(\varphi(x))\varphi'(x)$ holds for all $x\in U$. Due to \ref{assum:phi-psi-three}, it follows $\psi'(\varphi(f_t)) = 1/\varphi'(f_t) = 1/u_t$.
Together with \ref{assum:phi-psi-one}, we get
\begin{align*}
    \psi\big(\varphi(f_t) - \alpha_t v_t u_t^2 \|g_t\|^2\big) \geq f_t - \alpha_t v_t u_t \|g_t\|^2.
\end{align*}
%
Plugging back into \eqref{eqn:psi-phi-delta}, we have
\begin{align*}
    \delta_t &\leq f_t - \frac12 \alpha_t u_t^2 v_t^2 \|g_t\|^2 - \big(f_t - \alpha_t v_t u_t \|g_t\|^2\big)\\
    &= \frac{\alpha_t}{2} \|g_t\|^2 [2v_tu_t - u_t^2 v_t^2] \leq \frac{\alpha_t}{2} \|g_t\|^2,
\end{align*}
where the last step uses $1 -2u_tv_t + u_t^2 v_t^2 = (1-u_tv_t)^2 \geq 0$.
%


\paragraph{Using log and exponential.} Consider \eqref{eq:phi-psi-model} with $\psi = \exp$ and $\varphi = \log$, that is
\begin{align}\label{eqn:ngn-model-exp}
    f_{x}(y,s) = \exp \big(\log{f(x,s)} + \frac{1}{f(x,s)}\iprod{g}{y-x} \big).
\end{align}
The model-based update \eqref{eqn:update} has the necessary and sufficient optimality condition
\begin{align*}
    0 = \exp \big(\tfrac{1}{f(x_t,s_t)}\iprod{g_t}{y-x_t} \big) g_t + \frac{1}{\alpha_t}(y-x_t).
\end{align*}
Using the Ansatz $y=x_t - \gamma_t g_t$, we obtain
\begin{align}\label{eqn:exp-log-condition-gamma}
    \gamma_t = \alpha_t \exp \big(-\gamma_t\frac{\|g_t\|^2}{f(x_t,s_t)}\big).
\end{align}
Let $W_0$ be the principal branch of the Lambert--$W$ function, which is defined as the inverse function of $z\mapsto ze^z$. Multiplying \eqref{eqn:exp-log-condition-gamma} by $\frac{\|g_t\|^2}{f(x_t,s_t)}\exp \big(\gamma_t\frac{\|g_t\|^2}{f(x_t,s_t)}\big)$, we can see that \eqref{eqn:exp-log-condition-gamma} has the solution 
\begin{align}\label{eqn:lambertw-step-size}
    \gamma_t = \frac{f(x_t,s_t)}{\|g_t\|^2} W_0(\alpha_t \frac{\|g_t\|^2}{f(x_t,s_t)}) \quad \text{if } g_t \neq 0
    \text{ and else } \gamma_t=\alpha_t.
\end{align}
The solution lies on the principal branch as $\alpha_t \frac{\|g_t\|^2}{f(x_t,s_t)} > 0$.

Denote again $f_t = f(x_t,s_t)$. Further, we have $u_t = \frac{1}{f_t}$ and $\alpha_tu_tv_t = \gamma_t$. Hence, $v_t = \exp(\log f_t - \frac{\gamma_t}{f_t}\|g_t\|^2) = f_t \exp(- \frac{\gamma_t}{f_t}\|g_t\|^2) = \frac{f_t \gamma_t}{\alpha_t}$. We can plug this into \eqref{eqn:psi-phi-delta}, we get
\begin{align*}
    \delta_t &= f_t - \frac{\gamma_t^2}{2\alpha_t}\|g_t\|^2 - \exp(\log f_t - \frac{\gamma_t}{f_t}\|g_t\|^2) = f_t - \frac{\gamma_t^2}{2\alpha_t}\|g_t\|^2 - f_t\exp(- \frac{\gamma_t}{f_t}\|g_t\|^2) \\
    &= f_t - \frac{\gamma_t^2}{2\alpha_t}\|g_t\|^2 - \frac{f_t\gamma_t}{\alpha_t}.
\end{align*}

To the best of our knowledge, the adaptive step size \eqref{eqn:lambertw-step-size} is novel in the literature. In preliminary experiments, we observe that \eqref{eqn:lambertw-step-size} is also very stable with respect to the choice of $\alpha_t$. However, we do not further investigate the behavior of this new step size here, as it would distract from the focus of this paper. Analyzing \eqref{eqn:lambertw-step-size} might be interesting future work.


\paragraph{Datasets for logistic regression.}

We use the scaled versions of the datasets when available. We always use $20\%$ of the \texttt{LIBSVM} train split as validation set. Further, we drop the last batch if it is smaller than the batch size. This explains the difference between the sample size reported here and the one on the \texttt{LIBSVM} website.

We run 10 epochs for \texttt{dna} and 20 epochs for \texttt{vowel}.

\begin{table}[H]
    \centering
    \caption{Dataset information for multi-class logistic regression experiments.}
    \label{tab:logreg-datasets}
    \begin{tabular}{|c|cccc|}
         \hline
         Name & Dimension $d$ & Samples $n$ & Class labels $C$ & Batch size $b$ \\
         \hline
         \texttt{vowel} & 10 & 416 & 11 & 16 \\
         \texttt{dna} & 180 & 1600 & 3 & 16  \\
         \hline
    \end{tabular}
\end{table}

\section{Additional Plots and Experiments}

% Figure multiple D
\begin{figure}[H]
    \begin{subfigure}{0.49\columnwidth}
       \centering
        \includegraphics[width=0.99\linewidth]{plots/cifar10_resnet20/bound_last_allD.pdf}
        \caption{without warmup}
    \end{subfigure}
    %
    \begin{subfigure}{0.49\columnwidth}
       \centering
        \includegraphics[width=0.99\linewidth]{plots/cifar10_resnet20/warmup/bound_last_allD.pdf}
        \caption{with warmup}
    \end{subfigure}
    \caption{\texttt{ResNet20} on \texttt{CIFAR10}: Bound $\Omega_T$ from \cref{thm:last-iterate} for multiple values of $D \in [10^0, 10^2]$. The advantage of \SPS{} and \NGN{} over \SGD{} for large $\alpha$ is reflected in the bound independent of the value of $D$.}
    \label{fig:cifar10-multipleD}
\end{figure}
%
%
% Figure with train loss and bound, after only 9 epochs 
\begin{figure}[H]
    \begin{subfigure}{0.45\columnwidth}
       \centering
        \includegraphics[width=0.99\linewidth]{plots/cifar10_resnet20/stability_train_loss_epoch_9_bound.pdf}
        \caption{without warmup}
    \end{subfigure}
    %
    \begin{subfigure}{0.45\columnwidth}
       \centering
        \includegraphics[width=0.99\linewidth]{plots/cifar10_resnet20/warmup/stability_train_loss_epoch_9_bound.pdf}
        \caption{with warmup}
    \end{subfigure}
    \caption{\texttt{ResNet20} on \texttt{CIFAR10}: Same as \cref{fig:cifar10-loss-and-bound}, but computed after 10 instead of 20 epochs.}
    \label{fig:cifar10-loss-and-bound-ablate-T}
\end{figure}
%
%
% Figure with val loss and bound
\begin{figure}[H]
    \begin{subfigure}{0.45\columnwidth}
       \centering
        \includegraphics[width=0.99\linewidth]{plots/cifar10_resnet20/stability_val_loss_epoch_19_bound.pdf}
        \caption{without warmup}
    \end{subfigure}
    %
    \begin{subfigure}{0.45\columnwidth}
       \centering
        \includegraphics[width=0.99\linewidth]{plots/cifar10_resnet20/warmup/stability_val_loss_epoch_19_bound.pdf}
        \caption{with warmup}
    \end{subfigure}
    \caption{\texttt{ResNet20} on \texttt{CIFAR10}: Same as \cref{fig:cifar10-loss-and-bound}, but plotting validation loss instead of train loss.}
    \label{fig:cifar10-loss-and-bound-ablate-val}
\end{figure}
%

% DNA Figure 
\begin{figure}[H]
    \begin{subfigure}{0.45\columnwidth}
       \centering
        \includegraphics[width=0.99\linewidth]{plots/dna/stability_train_loss_epoch_9_bound.pdf}
    \end{subfigure}
    %
    \begin{subfigure}{0.45\columnwidth}
       \centering
        \includegraphics[width=0.99\linewidth]{plots/dna/deltas.pdf}
    \end{subfigure}
    \caption{Same as \cref{fig:vowel-bound-and-delta} but for logistic regression on \texttt{dna} dataset.}
    \label{fig:dna-bound-and-delta}
\end{figure}
%

%



\end{document}

% This document was modified from the file originally made available by
% Pat Langley and Andrea Danyluk for ICML-2K. This version was created
% by Iain Murray in 2018, and modified by Alexandre Bouchard in
% 2019 and 2021 and by Csaba Szepesvari, Gang Niu and Sivan Sabato in 2022.
% Modified again in 2023 and 2024 by Sivan Sabato and Jonathan Scarlett.
% Previous contributors include Dan Roy, Lise Getoor and Tobias
% Scheffer, which was slightly modified from the 2010 version by
% Thorsten Joachims & Johannes Fuernkranz, slightly modified from the
% 2009 version by Kiri Wagstaff and Sam Roweis's 2008 version, which is
% slightly modified from Prasad Tadepalli's 2007 version which is a
% lightly changed version of the previous year's version by Andrew
% Moore, which was in turn edited from those of Kristian Kersting and
% Codrina Lauth. Alex Smola contributed to the algorithmic style files.
